% Generated by build.py from article.tex and its section inputs.
\documentclass[11pt,a4paper]{article}

% BEGIN INCLUDED FILE: preamble.tex
\usepackage[T1]{fontenc}
\usepackage{lmodern}
\usepackage{amsmath,amssymb,amsthm,mathtools,mathrsfs}
\usepackage[a4paper,margin=26mm,headheight=15pt]{geometry}
\usepackage{microtype}
\usepackage{xcolor,booktabs,longtable,array,enumitem}
\usepackage{fancyhdr}
\usepackage{hyperref}
\usepackage{bookmark}
\usepackage{url}
\definecolor{ResearchBlue}{RGB}{28,66,91}
\hypersetup{colorlinks=true,linkcolor=ResearchBlue,citecolor=ResearchBlue,
 urlcolor=ResearchBlue,pdftitle={Bernoulli Cancellation, Harmonic Gaps, and Spectral Faces},
 pdfauthor={Research continuation prepared for Vladimir Reshetnikov},
 pdfsubject={Polylogarithms, polygamma tails, multiple zeta values, Stieltjes constants, and spectral defects}}
\urlstyle{same}
\setlength{\parskip}{3pt plus 1pt}
\setlength{\parindent}{14pt}
\setlength{\emergencystretch}{2em}
\allowdisplaybreaks[2]
\numberwithin{equation}{section}
\newtheorem{theorem}{Theorem}[section]
\newtheorem{proposition}[theorem]{Proposition}
\newtheorem{lemma}[theorem]{Lemma}
\newtheorem{corollary}[theorem]{Corollary}
\theoremstyle{definition}
\newtheorem{definition}[theorem]{Definition}
\theoremstyle{remark}
\newtheorem{remark}[theorem]{Remark}
\DeclareMathOperator{\CT}{CT}
\setlist[enumerate]{leftmargin=*,itemsep=5pt,topsep=5pt}
\setlist[itemize]{leftmargin=*,itemsep=3pt,topsep=4pt}
\pagestyle{fancy}
\fancyhf{}
\fancyhead[L]{\small\sffamily Bernoulli cancellation and spectral identities}
\fancyhead[R]{\small\sffamily ProveIt continuation}
\fancyfoot[C]{\thepage}
\renewcommand{\headrulewidth}{0.3pt}
\renewcommand{\footrulewidth}{0pt}

% END INCLUDED FILE: preamble.tex

\begin{document}
\hypersetup{pageanchor=false}
\begin{titlepage}
\vspace*{15mm}
{\sffamily\color{ResearchBlue}\large PROVEIT RESEARCH CONTINUATION\par}
\vspace{12mm}
{\LARGE\bfseries Bernoulli Cancellation,\par
Harmonic Gaps, and Spectral Faces\par}
\vspace{7mm}
{\large All-power trigamma moments, a completed reconstruction family,\par
and identities for Gamma and Stieltjes functions\par}
\vspace{13mm}
{\large Prepared for Vladimir Reshetnikov\par}
\vspace{3mm}
{11 October 2026\par}
\vspace{13mm}
\begin{abstract}
We continue several precise research questions in the ProveIt manuscript
and its five incoming reports of 11 October 2026. A single Bernoulli
polynomial identity cancels the finite-cutoff boundary terms for every
power of the trigamma tail. It gives an anchored analytic generator,
a finite multiple-zeta coordinate space at each power, a uniform
extension to all repeated even-order Hurwitz tails, and explicit
fifth- and sixth-power ordinary convergent sums. We determine all Laurent
coefficients on the complementary resonant lattice, including the exact
local convergence radius and Stieltjes normalization. A separate exact
certificate at the previously uncovered index completes the
smaller-order-ten optimal-prefix reconstruction family. For relative
multiple-Hurwitz interpolants we prove closure under finite decorations
and obtain several-marker identities involving log-Gamma products and
normalized Stieltjes primitives. Finally, an explicit subtraction on
every face of a Dirichlet simplex expresses all higher spectral defects
through lower one-mass derivatives; the full zero-jet correction is a
finite rational function and vanishes for mean-zero periodic weights.
Proofs, exact algebraic certificates, numerical diagnostics, a source
audit, and precise further questions are supplied. Statements of rank
refer to specified rational coordinate systems; no numerical period
independence or unproved Gaussian ladder identity is asserted.
\end{abstract}
\vfill
\noindent\textbf{Source baseline.}
\href{https://github.com/VladimirReshetnikov/ProveIt/tree/da2792fc9cc8ef19c6a314be6c65e1437fbdc66b}{ProveIt commit
\texttt{da2792fc9cc8ef19c6a314be6c65e1437fbdc66b}}.\par
\smallskip
\noindent\textbf{Status.} Research article with complete stated proof
chains and reproducible finite checks; prepared for subsequent repository
review and integration.
\end{titlepage}
\hypersetup{pageanchor=true}
\pagenumbering{roman}
\begingroup
\small
\setlength{\parskip}{0pt}
\tableofcontents
\endgroup
\clearpage
\pagenumbering{arabic}

% BEGIN INCLUDED FILE: sections/01_scope.tex
\section{Scope, results, and the source questions}
\label{sec:scope}

The source corpus is already extensive. It contains general harmonic
interpolation, centered tail moments, Gamma and Stieltjes derivative
ladders, finite spectral-subtraction algorithms, and a careful record
of rejected or still-conjectural period relations. A continuation should
therefore do more than rederive a familiar Euler sum or enlarge a
numerical search. We select explicit unresolved questions, identify
their inherited inputs, and supply complete new proof steps.

The baseline is commit
\texttt{da2792fc9cc8ef19c6a314be6c65e1437fbdc66b}, including the canonical
manuscript \cite{RepoMain}, the harmonic-gap continuation
\cite{gm:gaps-source}, and all five incoming archives dated 11 October
2026. In particular the quartic reconstruction report \cite{RepoQuartic}
already asks for fifth- and higher-power cancellation, the complementary
resonant parity, and completion of the order-ten reconstruction family.
The higher-defect report \cite{sd:HigherDefects} asks for explicit
simplex-face subtractions and their odd-weight extension. These are
concrete targets addressed here.

\subsection{The main new statements}

The first contribution is an all-power mechanism. If
\[
 A_n=\psi'(n+1)=\zeta(2)-H_n^{(2)},\qquad x_n=n+\tfrac12,
\]
the regular centered moments of $A_n^d$ have an anchored generating
function satisfying one triangular recurrence for every $d$. Its finite
boundary cancellation is the parity of
\[
 u^{2d}B_d\!\left(\tfrac12+\frac{E(u)}u\right),
 \qquad E(u)=\sum_{j\ge0}B_{2j}u^{2j}.
\]
This polynomial identity explains the lower-power examples and proves
the fixed finite-span conclusion without extrapolating from a finite
table. The resulting fifth-power coordinate system has eight formal
coordinates, and the sixth-power system has eleven. Exact rational
determinants certify the full ranks of the displayed coefficient rows.
An especially compact consequence is
\begin{equation}\label{intro:quintic}
\begin{split}
&\sum_{n=0}^{\infty}\Bigl[
 (195x_n+4056x_n^3+11568x_n^5+5376x_n^7)A_n^5
                  -5376x_n^2-9328\Bigr]\\
&\hspace{12mm}=-748\zeta(2)+7884\zeta(3)-5380\zeta(4)
                                    +171\zeta(6)-748.
\end{split}
\end{equation}
The entire bracket is one summand. Its series is ordinary and absolutely
convergent; splitting off its divergent polynomial pieces would change
the meaning. The proof is finite summation and elimination, not integer
relation recognition.

The same mechanism works for every repeated even-order Hurwitz tail
\(\zeta(2q,n+1)^d\), with an explicit finite span of weights at most
$2qd-1$ (Theorem~\ref{et:main}). A short consequence, with
$T_n=\psi^{(3)}(n+1)/6$, is
\[
 \sum_{n\ge0}\bigl((4x_n^3+x_n)T_n^3-3T_n^2\bigr)
                                      =-\frac94\zeta(8).
\]
The underlying cubic cutoff identity also gives a complex-order family
for $\Re p>3/2$. At the critical real order $p=3/2$, the ordinary
sum acquires the exact additional endpoint contribution $8$;
see Corollary~\ref{et:complex-cubic}.

The same spectral family has actual poles at $1-d-2r$. We give all their
Laurent coefficients as ordinary logarithmic sums plus a finite Hurwitz
and generalized-Stieltjes correction. Removing the pole gives a Taylor
series of exact radius two. This answers the normalization part of the
resonant question and supplies an analytic local generator, while leaving
open a uniform arithmetic reduction as the resonance index increases.

The second contribution completes a previously partial theorem. For
each even larger order $2b>10$, the first $2b$ regular centered moments
of $\zeta(10,n+1)\zeta(2b,n+1)$ have the optimal formal coordinate
rank $2b$. Two existing reports cover complementary ranges: a congruence
covers all but one residue class, and an eventual theorem covers
$b\ge1920$. Their only uncovered index is $b=962$. A five-by-five
matrix modulo $2^{20}$ proves nonvanishing there. We replay the inherited
finite polynomial identities and give an explicit rigorous large-order
comparison, so the completed family has no unchecked numerical case.
The additional parameter pair relative to the union of the two reports
is exactly $(10,1924)$.

The third contribution treats several transverse marks in relative
multiple-Hurwitz interpolation. Its first useful case is
\begin{equation}\label{intro:gamma}
\left.\partial_u\partial_v\mathcal H_{a,b}(s,u,v,t)\right|_{u=v=0}
=\frac12\mathcal E_{a,b}^{s,t}\!\left[
 (\log\Gamma(z+1)-\log\Gamma(y))^2
                  -2\int_y^{z+1}\gamma_1(x)\,dx\right].
\end{equation}
Here $\mathcal H$ is the prescribed relative quotient and $\mathcal E$
is an independently defined decorated double quotient; their full
definitions and ordinary-series realizations appear in
Section~\ref{gm:section}. The identity holds for arbitrary independent
spectral arguments and endpoints in the right half-plane. Its extension
uses a finite set-partition polynomial. Identical marked clusters give
individual ordered identities; different marks give their complete
permutation-symmetric sum. The analytic proof establishes decoration
closure before using the partition identity.

The fourth contribution resolves the face-subtraction problem for
higher spectral defects. For a polynomial weight, any number of
nonnegative masses, and every derivative order, we give an exact finite
sum of ordinary integrals over simplex faces. A $k$-vertex face first
appears at derivative order $k$, after the zero-jet part is separated.
That zero-jet part is an explicit rational function. It is absent for
even polynomial weights and, more generally, for every polynomial
weight multiplied by a mean-zero periodic sequence. The highest faces
at each derivative order require only the first one-mass derivative,
which for constant weight is an ordinary log-Gamma expression.
The subtraction also gives a closed identity at every face size:
\begin{equation}\label{intro:face}
 \int_{\Sigma_S}\mathcal H_S(B\mapsto B^2;t)
                  \frac{dt}{\prod_{i\in S}t_i}
 =(-1)^{|S|-1}(|S|-2)!
                 \sum_{\{i,j\}\subseteq S}(B_i-B_j)^2.
\end{equation}
The simplex normalization and face operator are specified in
Section~\ref{sd:sec:main}; the integral has no hidden finite part.

\subsection{Claim boundaries and relation to earlier work}

Several ingredients are classical and are used with attribution.
Hoffman's ordering of the minimum index gives products of zeta tails
as multiple zeta values \cite{HoffmanTails}. His symmetric-sum theorem
supplies the partition coefficients \cite{gm:Hoffman1992,gm:Hoffman2018}.
Ihara--Nakamura--Yamamoto supply the bare entire interpolation and its
harmonic product \cite{gm:INY}. The scalar Dirichlet mixture is a
classical Gamma-integral identity, already present in the source; the
literature on logarithmic Dirichlet averages includes
\cite{sd:Carlson1987,sd:Carlson1991}. These facts are inputs, not claimed
discoveries of the present article.

The new content is the all-power boundary mechanism and its concrete
consequences; the missing exact reconstruction certificate and completed
family; the local analytic realization of multiple decorations and
marks; and the explicit face decomposition with its rational correction
and all-order filtration. Novelty is assessed relative to the inspected
repository and the cited primary literature. An exhaustive priority
claim over the entire literature is not made.

The manuscript's Gaussian $S_6$ and revised $S_8$ period identities
remain open. The rejected historical $S_8$ vector is a different object
and retains its rejected status. No determinant of a formal coordinate
matrix proves arithmetic independence of evaluated periods. Likewise a
representation by polylogarithms of bounded depth is an upper bound in
that representation, not a proof of minimal depth. These distinctions
are part of the theorem statements, not qualifications added only after
the calculations.

\subsection{Conventions and proof architecture}

We use $H_n^{(r)}=\sum_{j=1}^n j^{-r}$, with $H_0^{(r)}=0$;
Bernoulli numbers have $B_1=-1/2$; and multiple zeta values are ordered
as
\[
 \zeta(s_1,\ldots,s_\ell)
 =\sum_{n_1>\cdots>n_\ell\ge1}
                         n_1^{-s_1}\cdots n_\ell^{-s_\ell}.
\]
The prime on a zeta function denotes differentiation in its spectral
argument unless explicitly stated otherwise. The generalized Stieltjes
constants satisfy
\[
 \zeta(1+\varepsilon,a)=\frac1\varepsilon+
       \sum_{q\ge0}\frac{(-1)^q\gamma_q(a)}{q!}\varepsilon^q.
\]
Section-specific letters are reset explicitly: $F_d(s)$ denotes a tail
spectrum, whereas $F_j(B)$ in the spectral-face section denotes a
one-mass derivative. Every antiderivative has its endpoint fixed.

The arguments use three distinct kinds of evidence. Finite identities
are proved algebraically and supplied with rational or modular replay.
Analytic statements include convergence domains, continuation from an
open domain, and endpoint normalization. High-precision evaluations
then check independent representations and signs. Numerical agreement
is not used as a substitute for either of the first two steps.
Section~\ref{sec:audit} records what was audited and what the artifact
checks establish; Section~\ref{sec:research} lists exact next targets.

% END INCLUDED FILE: sections/01_scope.tex


% BEGIN INCLUDED FILE: sections/02_all_powers.tex
\section{A Bernoulli cancellation law for every trigamma power}
\label{tp:sec:allpowers}

The quartic report asks for a cancellation mechanism at the fifth and
higher powers, a fixed arithmetic coordinate space at each power, and a
generator with its constants specified \cite[Research question R4]{RepoQuartic}.
We give one recurrence for all powers.  It explains the earlier quadratic,
cubic, and quartic cancellations and yields new fifth- and sixth-power
formulas.  The essential ingredient is a Bernoulli polynomial, rather than
a new summation theorem for each successive power.

\subsection{Regular moments as ordinary convergent sums}

Throughout this section let
\begin{equation}\label{tp:eq:tails}
 A_n=\psi'(n+1)=\zeta(2)-H_n^{(2)},\qquad x_n=n+\tfrac12,
 \qquad n\geq0.
\end{equation}
Bernoulli numbers have the convention $B_1=-1/2$.  For $d\geq1$ define
\begin{equation}\label{tp:eq:h}
 h_{d,j}=\sum_{j_1+\cdots+j_d=j}\prod_{\nu=1}^d B_{2j_\nu}(1/2),
 \qquad j_\nu\geq0.
\end{equation}
In particular $h_{d,0}=1$ and $h_{d,1}=-d/12$.  The centered
Hurwitz expansion \cite[\S25.11]{DLMF} gives
\begin{equation}\label{tp:eq:centeredasy}
 A_n^d\sim\sum_{j\geq0}h_{d,j}x_n^{-d-2j}.
\end{equation}
This is an asymptotic expansion; it is not a convergent Laurent series
at a fixed finite $n$.

\begin{proposition}[Ordinary realization of the regular parity]
\label{tp:prop:ordinary}
The function
\[
 F_d(s)=\sum_{n\geq0}A_n^d x_n^{-s},\qquad \Re s>1-d,
\]
has a meromorphic continuation to $\mathbb C$.  Its poles are exactly
$1-d-2j$, $j\geq0$, and their residues are $h_{d,j}$.  For every
$m\geq0$ with $m\equiv d\pmod2$, the value $M_{d,m}=F_d(-m)$ is regular
and is the absolutely convergent sum
\begin{equation}\label{tp:eq:ordinary}
 M_{d,m}=\sum_{n\geq0}\left(
 x_n^m A_n^d-
 \sum_{0\leq j\leq(m-d)/2}h_{d,j}x_n^{m-d-2j}\right).
\end{equation}
An inner sum with negative upper limit is empty.
\end{proposition}

\begin{proof}
Subtract a finite number of terms of \eqref{tp:eq:centeredasy} and add
the corresponding functions $h_{d,j}\zeta(s+d+2j,1/2)$.  Increasing
the subtraction order gives normally convergent remainders in
successively larger half-planes.  The meromorphic continuations agree
on overlaps, since they agree with the defining series in its initial
half-plane.  The signs of $B_{2j}(1/2)$ are $(-1)^j$, including $j=0$.
Every summand in \eqref{tp:eq:h} therefore has sign $(-1)^j$; hence
$h_{d,j}\ne0$ and the listed poles are actual.

At $s=-m$ the terms subtracted in \eqref{tp:eq:ordinary} contribute
$h_{d,j}\zeta(d+2j-m,1/2)=0$.  Indeed their Hurwitz arguments are
nonpositive even integers, and $\zeta(-2q,1/2)=0$ for $q\geq0$.
The next summand is $O(x_n^{-2})$, or decays faster when no subtraction
is needed.  This proves both the ordinary convergence and the equality.
\end{proof}

Thus the word ``moment'' below always refers to the explicit convergent
expression \eqref{tp:eq:ordinary}.  The meromorphic family is a proof
and differentiation tool; no assignment of a value to an unspecified
divergent sum is required.

\subsection{Finite-cutoff coordinates}

For a finite-cutoff asymptotic expansion, let $\CT_N$ denote the coefficient
of $N^0(\log N)^0$.  Put
\begin{equation}\label{tp:eq:TC}
 \begin{split}
 T_{e,q}&=\CT_N\sum_{n=1}^N n^q A_n^e,\qquad
 b_e(k)=[u^k]\mathcal A(u)^e,\\
 \mathcal A(u)&=\sum_{a\geq1}B_{a-1}u^a,\\
 C_{d,k}&=b_d(k)+\sum_{j=1}^d\binom dj T_{d-j,k-2j}.
 \end{split}
\end{equation}
Here $b_0(k)=\delta_{k,0}$, and a coefficient with negative index is
zero.  The series $\mathcal A$ is formal.  These constant coefficients
exist: repeated finite summation by parts reduces their existence to
ordinary power sums, harmonic numbers, and convergent multiple zeta
values.  A constructive reduction is given below.

The exact recurrence
$A_{n-1}=A_n+n^{-2}$ shows that $C_{d,k}$ is the constant coefficient
of
\[
 \sum_{n=1}^N n^k(A_{n-1}^d-A_n^d)+N^kA_N^d.
\]
Define the centered Faulhaber polynomial
\begin{equation}\label{tp:eq:Faulhaber}
 P_m(N)=\sum_{n=0}^{N-1}(n+1/2)^m
 =\frac{B_{m+1}(N+1/2)-B_{m+1}(1/2)}{m+1}.
\end{equation}
Finite summation by parts gives
\begin{equation}\label{tp:eq:finiteSBP}
 \sum_{n=0}^{N-1}x_n^m A_n^d
 =P_m(N)A_N^d+
   \sum_{n=1}^N P_m(n)(A_{n-1}^d-A_n^d).
\end{equation}
Every subtracted even-power sum in \eqref{tp:eq:ordinary} is an odd
polynomial in $N$, so it has zero constant coefficient.  Consequently,
if $P_m(N)=\sum_{k\geq1}p_{m,k}N^k$, then
\begin{equation}\label{tp:eq:momentC}
 M_{d,m}=\sum_{k\geq1}p_{m,k}C_{d,k}.
\end{equation}

The parity in \eqref{tp:eq:Faulhaber} is useful.  For even $m$ the
polynomial is odd; for odd $m$ it is even with zero constant term.
Formally define
\begin{equation}\label{tp:eq:GH}
 G_d(t)=\sum_{\substack{m\geq0\\m\equiv d\ (2)}}M_{d,m}\frac{t^m}{m!},
 \qquad H_d(t)=2\sinh(t/2)G_d(t).
\end{equation}
The finite polynomial identity \eqref{tp:eq:momentC} is equivalent to
\begin{equation}\label{tp:eq:HC}
 H_d(t)=\sum_{\substack{k\geq1\\k\equiv d+1\ (2)}}C_{d,k}\frac{t^k}{k!}.
\end{equation}
For example, summing \eqref{tp:eq:Faulhaber} over even $m$ gives
$\sinh(Nt)/(2\sinh(t/2))$; summing over odd $m$ gives
$(\cosh(Nt)-1)/(2\sinh(t/2))$.  These two identities prove
\eqref{tp:eq:HC} coefficient by coefficient without assuming convergence
of a moment generating series in advance.

\subsection{The universal boundary cancellation}

\begin{lemma}[A finite Bernoulli difference identity]
\label{tp:lem:EM}
For every integer $d\geq1$ and indeterminates $a,h$,
\begin{equation}\label{tp:eq:EM}
 \begin{split}
 (a+h)^d-a^d={}&\frac{dh}{2}\bigl((a+h)^{d-1}+a^{d-1}\bigr)\\
 &-\sum_{j=1}^{\lfloor(d-1)/2\rfloor}
      \binom d{2j}B_{2j}h^{2j}
      \bigl((a+h)^{d-2j}-a^{d-2j}\bigr).
 \end{split}
\end{equation}
Moreover, whenever $k\equiv d+1\pmod2$,
\begin{equation}\label{tp:eq:boundary}
 b_d(k)+\frac d2 b_{d-1}(k-2)
 +\sum_{j=1}^{\lfloor d/2\rfloor}
      \binom d{2j}B_{2j}b_{d-2j}(k-4j)=0.
\end{equation}
\end{lemma}

\begin{proof}
The first identity is the terminating Euler--Maclaurin identity for a
polynomial on one interval.  It also follows by comparing coefficients
in the Bernoulli generating function, so its use here is purely finite
algebra.  A term with $d-2j=0$ would be a difference of constants and is
omitted.

For the second identity write
\[
 E(u)=\sum_{j\geq0}B_{2j}u^{2j},\qquad
 \mathcal A(u)=uE(u)-u^2/2.
\]
The expression on the left of \eqref{tp:eq:boundary} is the coefficient
of $u^k$ in
\begin{equation}\label{tp:eq:BernoulliMechanism}
 \begin{split}
 &\mathcal A(u)^d+\frac d2u^2\mathcal A(u)^{d-1}
 +\sum_{j\geq1}\binom d{2j}B_{2j}u^{4j}\mathcal A(u)^{d-2j}\\
 &\hspace{32mm}=u^{2d}B_d\left(\frac12+\frac{E(u)}u\right).
 \end{split}
\end{equation}
The Bernoulli reflection law $B_d(1-z)=(-1)^dB_d(z)$ implies that
the right side has parity $d$ as a formal series in $u$.  Its
coefficient at the opposite parity is zero.  Although the argument of
$B_d$ contains $u^{-1}$, the displayed product has no negative powers:
the polynomial has degree $d$.  This establishes the cancellation.
\end{proof}

Identity \eqref{tp:eq:BernoulliMechanism} supplies a single reason for
all the boundary terms to cancel.  The separate quartic cancellation
is one low-degree consequence.

\begin{theorem}[All-power triangular recurrence]
\label{tp:thm:recurrence}
Set $G_0=0$.  For $d\geq1$ and $k\geq2$ with
$k\equiv d+1\pmod2$,
\begin{equation}\label{tp:eq:recurrence}
 \boxed{\begin{aligned}
 C_{d,k}={}&d(k-2)![t^{k-2}]\bigl(\cosh(t/2)G_{d-1}(t)\bigr)
       -\frac d2\delta_{k,2}\zeta(2)^{d-1}\\
 &-\sum_{j=1}^{\lfloor(d-1)/2\rfloor}
       \binom d{2j}B_{2j}C_{d-2j,k-4j}.
 \end{aligned}}
\end{equation}
Every coordinate on the right has smaller tail power.
\end{theorem}

\begin{proof}
Apply \eqref{tp:eq:EM} with $a=A_n$ and $h=n^{-2}$, multiply by
$n^k$, and sum only up to a finite $N$.  For $m=k-2\geq0$, changing
the index in one of the two symmetric sums gives the exact relation
\begin{align*}
 &\sum_{n=1}^N n^m(A_{n-1}^{d-1}+A_n^{d-1})\\
 &\quad=\sum_{n=0}^{N-1}\bigl((n+1)^m+n^m\bigr)A_n^{d-1}
       -\delta_{m,0}\zeta(2)^{d-1}+N^mA_N^{d-1}.
\end{align*}
The polynomial $(x+1/2)^m+(x-1/2)^m$ contains only powers of
parity $m\equiv d-1\pmod2$.  Taking its constant coefficient therefore
uses exactly the regular moments of power $d-1$.  Its contribution is
$2m![t^m](\cosh(t/2)G_{d-1})$.  For $d=1$, these polynomial sums have
zero constant coefficient, which is the convention $G_0=0$.

The differences of the lower powers in \eqref{tp:eq:EM} give
$C_{d-2j,k-4j}-b_{d-2j}(k-4j)$.  The remaining boundary terms are
precisely \eqref{tp:eq:boundary}; their sum is zero.  Taking constant
coefficients after these finite identities have been combined proves
\eqref{tp:eq:recurrence}.  In particular the exceptional $k=2$ term
comes from the missing $n=0$ endpoint and cannot be discarded.
\end{proof}

\subsection{Explicit anchors and a convergent polylogarithm generator}

Write $\boldsymbol a\vDash e$ for a composition of $e$.  For $e\geq1$
and $h\leq1$, excluding $(e,h)=(1,1)$, define the finite MZV combination
\begin{equation}\label{tp:eq:anchor}
 \mathcal Z_{e,h}=
 \sum_{(a_1,\ldots,a_\ell)\vDash e}
 \frac{e!}{a_1!\cdots a_\ell!}
 \zeta(2a_1,\ldots,2a_{\ell-1},2a_\ell-h).
\end{equation}
All the displayed multiple zeta values converge.  Ordering the indices
by their minimum, as in Hoffman's tail-product method \cite{HoffmanTails},
gives
\begin{equation}\label{tp:eq:anchorseries}
 C_{e,h}=\mathcal Z_{e,h}
 =\sum_{n\geq1}n^h(A_{n-1}^e-A_n^e),\qquad h\leq1.
\end{equation}
The exceptional pair is divergent and is never needed.  For $h=0$,
telescoping yields $\mathcal Z_{e,0}=\zeta(2)^e$.  Formula
\eqref{tp:eq:anchor} is classical index ordering, not a new claim about
unweighted tail products.

Let $I^q$ denote an anchored $q$-fold primitive:
\begin{equation}\label{tp:eq:primitive}
 (I^qf)(t)=\frac1{(q-1)!}\int_0^t(t-u)^{q-1}f(u)\,du,
 \qquad q\geq1.
\end{equation}
For $d\geq1$ put $c_{d,j}=\binom d{2j}B_{2j}$ and let
\begin{equation}\label{tp:eq:anchorpoly}
 \begin{split}
 R_d(t)={}&\mathbf1_{2\mid d}\,\mathcal Z_{d,1}t
       -\mathbf1_{2\nmid d}\,\frac d4\zeta(2)^{d-1}t^2\\
 &-\sum_{j=1}^{\lfloor(d-1)/2\rfloor}c_{d,j}
       \sum_{\substack{2\leq k\leq4j\\k\equiv d+1\ (2)}}
         \mathcal Z_{d-2j,k-4j}\frac{t^k}{k!}.
 \end{split}
\end{equation}
For $d=1$ the first term is absent; it is not an invocation of the
undefined anchor $\mathcal Z_{1,1}$.

\begin{theorem}[An anchored generator at every power]
\label{tp:thm:generator}
The series in \eqref{tp:eq:GH} converge for $|t|<2\pi$.  They are
determined recursively by $G_0=H_0=0$ and
\begin{equation}\label{tp:eq:generator}
 \boxed{\begin{aligned}
 H_d(t)={}&R_d(t)+dI^2\bigl(\cosh(t/2)G_{d-1}(t)\bigr)\\
 &-\sum_{j=1}^{\lfloor(d-1)/2\rfloor}c_{d,j}I^{4j}H_{d-2j}(t),
 \qquad G_d(t)=\frac{H_d(t)}{2\sinh(t/2)}.
 \end{aligned}}
\end{equation}
The apparent singularity at zero is removable.  At each fixed $d$,
all moments $M_{d,m}$ lie in one explicitly specified finite rational
span of convergent MZVs and $1$, independent of $m$.
One uniform, possibly redundant choice is the span of $1$ and all
positive-index MZVs of weight at most $2d-1$ with first index at least two.

Moreover $H_d$ has a finite expression in polynomials in $t$, multiple
zeta constants, and one-variable multiple polylogarithms at $e^{-t}$
of functional depth at most $d-1$.  This is an upper bound in a
specified representation, not a minimal-depth or period-independence claim.
\end{theorem}

\begin{proof}
Sum \eqref{tp:eq:recurrence} with the factors $t^k/k!$.  The positive
lower indices give $I^{4j}H_{d-2j}$.  The nonpositive lower indices
give exactly the finite polynomial in \eqref{tp:eq:anchorpoly}, and
the $k=1$ coefficient at even $d$ is \eqref{tp:eq:anchorseries}.  Thus
\eqref{tp:eq:generator} holds first as a formal power-series identity.

Inductively the numerator is holomorphic for $|t|<2\pi$ and has
parity $d+1$.  It is $O(t)$ at even $d$ and $O(t^2)$ at odd $d$.
The denominator has only a simple zero at zero inside this disk.
Division therefore preserves holomorphy and gives parity $d$.
This proves convergence of the formal moment series and the asserted
value of its coefficients.  At fixed $d$ only finitely many anchors
from powers at most $d$ have been introduced; their weights are at
most $2d-1$.  Taylor coefficients of all remaining elementary factors
and integrations are rational.  This proves the finite-span assertion.

For the functional representation start with real $t>0$ and put
$z=e^{-t}$. The assembled expression will then denote the analytic
continuation supplied by the anchored generator on $|t|<2\pi$;
separated terms need not be holomorphic at zero. The factor that
enters the new integration step is
\[
 \cosh(t/2)G_{d-1}(t)
 =\left(\frac12+\frac z{1-z}\right)H_{d-1}(t).
\]
Use the convention
\[
 L_{a_1,\ldots,a_\ell}(t)
 =\sum_{n_1>\cdots>n_\ell\geq1}
       \frac{e^{-tn_1}}{n_1^{a_1}\cdots n_\ell^{a_\ell}}.
\]
The empty word represents $1$. We maintain the stronger induction
invariant that each nonempty polylogarithm in $H_d$ has first index at
least two. Ordinary integrations raise that first index. For the new
operation, an explicit double primitive of
$t^m zL_{\boldsymbol a}(t)/(1-z)$ is
\[
 P_{m,\boldsymbol a}(t)=\sum_{j=0}^m(j+1)\frac{m!}{(m-j)!}
                t^{m-j}L_{j+2,\boldsymbol a}(t).
\]
Indeed $L_{1,\boldsymbol a}'=-zL_{\boldsymbol a}/(1-z)$ and
$L_{a+1,\boldsymbol b}'=-L_{a,\boldsymbol b}$. This raises depth by
at most one and preserves the leading-index invariant. Its value at
zero is a convergent MZV, and its first derivative is
\[
 P_{m,\boldsymbol a}'(t)=-\sum_{j=0}^m\frac{m!}{(m-j)!}
                t^{m-j}L_{j+1,\boldsymbol a}(t).
\]
For $m>0$ the endpoint is $-m!\zeta(m+1,\boldsymbol a)$. For $m=0$
the only potential divergence has exactly one leading index one.

For an admissible word $\boldsymbol a=(a_1,\ldots,a_\ell)$ define
\begin{align*}
 K(\boldsymbol a)={}&-\sum_{i=1}^{\ell}
 \zeta(a_1,\ldots,a_i,1,a_{i+1},\ldots,a_\ell)\\
 &-\sum_{i=1}^{\ell}
 \zeta(a_1,\ldots,a_i+1,\ldots,a_\ell),
 \qquad K(\varnothing)=0.
\end{align*}
All these values converge, and
\[
 L_{1,\boldsymbol a}(t)
       =\zeta(\boldsymbol a)L_1(t)+K(\boldsymbol a)+o(1).
\]
To prove this, subtract $\zeta(\boldsymbol a)L_1(t)$ termwise.
The remaining series is
$\sum_{n\ge1}e^{-tn}(H_{n-1}(\boldsymbol a)-\zeta(\boldsymbol a))/n$,
where $H_{n-1}(\boldsymbol a)$ is the finite multiple harmonic sum.
It is absolutely convergent at $t=0$ because $a_1\ge2$. Ordering its
extra index among the zeta indices gives exactly the two sums defining
$K$. Dominated convergence proves the displayed endpoint formula.

In the complete double primitive the coefficient of the potentially
divergent $L_1(t)$ is $-H_{d-1}(0)=0$. Every remaining endpoint is
therefore an explicit convergent MZV. Subtracting the value and first
derivative at zero anchors the primitive. The other integrations in
\eqref{tp:eq:generator} do not increase depth. This proves the bound
$d-1$ and fixes every integration constant without assigning values
to individual divergent multiple zeta series.
\end{proof}

For completeness, if $F_j'=-F_{j-1}$, an unanchored $q$-fold primitive
of $t^mF_0(t)$ is
\begin{equation}\label{tp:eq:finiteprimitive}
 (-1)^q\sum_{j=0}^m\binom{q+j-1}{j}
       \frac{m!}{(m-j)!}t^{m-j}F_{q+j}(t).
\end{equation}
Subtracting its Taylor polynomial of degree $q-1$ at zero supplies
the anchor, with grouped limits when necessary.  Thus
\eqref{tp:eq:generator} is also a finite algorithm for antiderivatives
and multiple-polylogarithm formulas, rather than a recurrence with
unspecified constants.

\subsection{A finite-cutoff implementation independent of the new recurrence}

The supplied verification program does not use \eqref{tp:eq:generator}
to construct the comparison moments.  If
$F_q(N)=\sum_{n=1}^N n^q=\sum_a f_aN^a$, ordinary finite summation
by parts instead gives
\begin{equation}\label{tp:eq:independentT}
 T_{e,q}=\sum_a f_ab_e(a)+
 \sum_{j=1}^e\binom ej
 \left(\sum_a f_aT_{e-j,a-2j}-T_{e-j,q-2j}\right),\qquad q\geq0.
\end{equation}
For $q<0$, products of strictly greater tail indices give
\[
 T_{e,q}=\sum_{\boldsymbol a\vDash e}
       \frac{e!}{a_1!\cdots a_\ell!}
       \zeta(2a_1,\ldots,2a_\ell,-q),\qquad e\geq1.
\]
The remaining base cases are $T_{0,q}=0$ for $q\geq0$,
$T_{0,-1}=\gamma$, and $T_{0,-p}=\zeta(p)$ for $p\geq2$.
The power $e$ decreases at every recursive step.  This proves existence
of the constant coefficients used above and supplies a terminating
independent exact verification route.

\subsection{The same cancellation for every even Hurwitz order}
\label{et:section}

Fix $q\geq1$ and put $p=2q$. In this subsection the symbols
$M,C,G,H$ refer to the fixed tail
$T_n=\zeta(p,n+1)=\psi^{(p-1)}(n+1)/(p-1)!$, with
$x_n=n+1/2$. Thus $T_{n-1}=T_n+n^{-p}$. Define
\[
 \eta_0=\frac1{p-1},\qquad
 \eta_j=\frac{(p)_{2j-1}}{(2j)!}B_{2j}(1/2)\quad(j\geq1),
 \qquad h_{d,j}=\sum_{j_1+\cdots+j_d=j}\prod_{\nu=1}^d\eta_{j_\nu}.
\]
The centered Hurwitz asymptotics give
$T_n^d\sim\sum_{j\geq0}h_{d,j}x_n^{-(p-1)d-2j}$.
For $m\geq0$, $m\equiv d\pmod2$, the regular continued moment is
the ordinary absolutely convergent sum
\begin{equation}\label{et:ordinary}
 M_{d,m}=\sum_{n\geq0}\left\{x_n^mT_n^d-
 \sum_{0\leq j\leq(m-(p-1)d)/2}
              h_{d,j}x_n^{m-(p-1)d-2j}\right\}.
\end{equation}
Indeed the next term is $O(x_n^{-2})$, while every removed term
has continued sum $\zeta(-2r,1/2)=0$. The meromorphic family
$\sum T_n^dx_n^{-s}$ has its poles at $1-(p-1)d-2j$;
the same-sign convolution of the $\eta_j$ shows that their residues
$h_{d,j}$ are nonzero.

For every integer $k$ put
\[
 C_{d,k}=\CT_N\left\{\sum_{n=1}^N n^k(T_{n-1}^d-T_n^d)
                                  +N^kT_N^d\right\},\qquad
 G_d(t)=\sum_{m\equiv d\ (2)}M_{d,m}\frac{t^m}{m!}.
\]
Subtracting finitely many asymptotic terms proves existence of each
cutoff constant. The latter sum has $m\geq0$.
Set $H_d=2\sinh(t/2)G_d$ and
$G_0=H_0=0$. The centered Faulhaber identity used above gives
\begin{equation}\label{et:HC}
 H_d(t)=\sum_{\substack{k\geq1\\k\equiv d+1\ (2)}}C_{d,k}\frac{t^k}{k!}.
\end{equation}
For $e\geq1$ and $h<p$, except $(e,h)=(1,p-1)$, the needed anchors are
\begin{equation}\label{et:anchors}
 \mathcal Z^{(q)}_{e,h}=
 \sum_{(a_1,\ldots,a_\ell)\vDash e}
 \frac{e!}{a_1!\cdots a_\ell!}
 \zeta(pa_1,\ldots,pa_{\ell-1},pa_\ell-h)=C_{e,h}.
\end{equation}
Here $\vDash$ denotes a composition. Ordering the $e$ tail indices
by their minimum proves this formula: the summand is
$n^h(T_{n-1}^e-T_n^e)=O(n^{h-(p-1)e-1})$ and the boundary term
tends to zero. The excluded pair never has the required parity.

Write $c_{d,j}=\binom d{2j}B_{2j}$ and define the finite polynomial
\begin{align}\label{et:R}
 R_d(t)={}&\sum_{\substack{1\leq k<p\\k\equiv d+1\ (2)}}
       \mathcal Z^{(q)}_{d,k}\frac{t^k}{k!}
 -\mathbf1_{2\nmid d}\frac d2\zeta(p)^{d-1}\frac{t^p}{p!}\notag\\
 &-\sum_{j=1}^{\lfloor(d-1)/2\rfloor}c_{d,j}
   \sum_{\substack{p\leq k\leq2pj\\k\equiv d+1\ (2)}}
       \mathcal Z^{(q)}_{d-2j,k-2pj}\frac{t^k}{k!}.
\end{align}

\begin{theorem}[Every power of every even-order tail]\label{et:main}
For $k\geq p$ with $k\equiv d+1\pmod2$,
\begin{equation}\label{et:recurrence}
 \begin{split}
 C_{d,k}={}&d(k-p)![t^{k-p}]\bigl(\cosh(t/2)G_{d-1}(t)\bigr)
       -\frac d2\delta_{k,p}\zeta(p)^{d-1}\\
 &-\sum_{j=1}^{\lfloor(d-1)/2\rfloor}c_{d,j}C_{d-2j,k-2pj}.
 \end{split}
\end{equation}
With the anchored primitives of \eqref{tp:eq:primitive}, its generator is
\begin{equation}\label{et:generator}
 H_d=R_d+dI^p\bigl(\cosh(t/2)G_{d-1}\bigr)
        -\sum_{j=1}^{\lfloor(d-1)/2\rfloor}c_{d,j}I^{2pj}H_{d-2j},
 \qquad G_d=\frac{H_d}{2\sinh(t/2)}.
\end{equation}
Both functions are holomorphic for $|t|<2\pi$. At fixed $(q,d)$,
all the moments belong to a fixed finite rational span of $1$ and
convergent MZVs of weights at most $2qd-1$. The generator $H_d$ admits
a finite multiple-polylogarithm representation of functional depth at
most $d-1$, with all constants fixed by the anchored primitives.
\end{theorem}

\begin{proof}
The uncentered formal asymptotic series is
\[
 \mathcal A_q(u)=u^{p-1}E_q(u)-\tfrac12u^p,
 \qquad E_q(u)=\frac1{p-1}
    +\sum_{j\geq1}\frac{(p)_{2j-1}B_{2j}}{(2j)!}u^{2j}.
\]
Its boundary combination satisfies the exact formal identity
\begin{equation}\label{et:boundary}
 \mathcal A_q^d+\frac d2u^p\mathcal A_q^{d-1}
 +\sum_{j=1}^{\lfloor d/2\rfloor}c_{d,j}u^{2pj}\mathcal A_q^{d-2j}
 =u^{pd}B_d\left(\frac12+\frac{E_q(u)}u\right).
\end{equation}
Since $E_q$ is even, Bernoulli reflection makes the right side have
parity $d$. Thus the required opposite-parity coefficients vanish.
Apply the finite polynomial identity \eqref{tp:eq:EM} with
$a=T_n$, $h=n^{-p}$ before taking constant coefficients. For
$r=k-p\geq0$, the symmetric term has the exact cutoff shift
\[
 \sum_{n=1}^N n^r(T_{n-1}^{d-1}+T_n^{d-1})
 =\sum_{n=0}^{N-1}\bigl((n+1)^r+n^r\bigr)T_n^{d-1}
   +N^rT_N^{d-1}-\delta_{r,0}\zeta(p)^{d-1}.
\]
The centered bracket has regular parity and gives
$2r![t^r](\cosh(t/2)G_{d-1})$. Equation~\eqref{et:boundary}
cancels the remaining boundaries, proving \eqref{et:recurrence}.
For $d=1$ the polynomial sums have zero constant coefficient, in
agreement with $G_0=0$.

Summing the recurrence gives \eqref{et:generator}: the low positive
indices give the first term of $R_d$, and the nonpositive lower
indices give its last term. Induction gives holomorphy on the stated
disk; $H_d$ has parity $d+1$ and vanishes at zero, so the division is
removable. Each low anchor has weight $pe-h\leq pd-1$; each last-term
anchor has weight $p(d-2j)-(k-2pj)=pd-k\leq pd-p$.
All other Taylor operations are rational. Finally
$\cosh(t/2)G_{d-1}=(1/2+z/(1-z))H_{d-1}$, $z=e^{-t}$;
the same primitive construction as in Theorem~\ref{tp:thm:generator}
raises functional depth by at most one. Increasing the number of
anchored integrations from $2$ to $p$ adds no depth.
\end{proof}

This extends the inspected two-tail and trigamma-power results to every
repeated even order. It does not yet address unequal orders among
three or more factors. No priority claim beyond the inspected source
corpus is intended.

\begin{corollary}[An ordinary cubic identity at every even order]
\label{et:cubic}
For every $q\geq1$, with $T_n=\zeta(2q,n+1)$,
\begin{equation}\label{et:cubic-formula}
 \boxed{\sum_{n\geq0}\left\{
       \bigl((n+1)^{2q}-n^{2q}\bigr)T_n^3-3T_n^2\right\}
       =-\frac{3\zeta(2q)^2+\zeta(4q)}2.}
\end{equation}
In particular, for $x_n=n+1/2$ and $T_n=\psi^{(3)}(n+1)/6$,
\begin{equation}\label{et:psi3}
 \boxed{\sum_{n\geq0}
       \bigl((4x_n^3+x_n)T_n^3-3T_n^2\bigr)
                      =-\frac94\zeta(8).}
\end{equation}
Both identities are absolutely convergent without subtraction.
\end{corollary}

\begin{proof}
Put $p=2q$. The cubic case of \eqref{tp:eq:EM}, followed by finite
summation by parts, gives the exact cutoff identity
\begin{align*}
 &\sum_{n=0}^{N-1}\left\{
       \bigl((n+1)^p-n^p\bigr)T_n^3-3T_n^2\right\}\\
 &\qquad=N^pT_N^3+\frac32T_N^2
       -\frac32\zeta(p)^2-\frac12H_N^{(2p)}.
\end{align*}
Here the cubic difference is
$n^p(T_{n-1}^3-T_n^3)
 =\tfrac32(T_{n-1}^2+T_n^2)-\tfrac12n^{-2p}$.
Since $T_n=O((n+1)^{1-p})$, the two boundary terms tend to zero
and each series on the left converges absolutely for $p\geq2$.
Letting $N\to\infty$ proves \eqref{et:cubic-formula}.
For $p=4$, use $(n+1)^4-n^4=4x_n^3+x_n$ and
$\zeta(4)^2=7\zeta(8)/6$ to obtain \eqref{et:psi3}.
\end{proof}

\begin{corollary}[Complex orders and a critical cubic value]
\label{et:complex-cubic}
Put $T_n(p)=\zeta(p,n+1)$, using the real logarithm of positive
integers in all powers and setting $0^p=0$ in the domain below.
For every complex $p$ with $\Re p>3/2$,
\begin{equation}\label{et:complex-cubic-formula}
 \sum_{n\geq0}\left\{
    \bigl((n+1)^p-n^p\bigr)T_n(p)^3-3T_n(p)^2\right\}
       =-\frac{3\zeta(p)^2+\zeta(2p)}2.
\end{equation}
At the real critical order $p=3/2$, the combined series also converges
absolutely, but its value has an additional boundary term:
\begin{equation}\label{et:critical-cubic}
 \boxed{\begin{aligned}
 &\sum_{n\geq0}\left\{
    \bigl((n+1)^{3/2}-n^{3/2}\bigr)T_n(3/2)^3-3T_n(3/2)^2\right\}
 \\[-1mm]
 &\qquad=8-\frac{3\zeta(3/2)^2+\zeta(3)}2.
 \end{aligned}}
\end{equation}
\end{corollary}

\begin{proof}
The preceding cutoff identity remains exact for $\Re p>1$.
For fixed $p$ in this half-plane,
$T_N(p)=N^{1-p}/(p-1)+O(N^{-\Re p})$.
If $\Re p>3/2$, each series on its left converges absolutely and
$N^pT_N(p)^3=O(N^{3-2\Re p})\to0$, proving
\eqref{et:complex-cubic-formula}.

At $p=3/2$, put $x=n+1/2$. The centered expansions are
\[
 T_n(3/2)=2x^{-1/2}-\frac1{16}x^{-5/2}+O(x^{-9/2}),\qquad
 (n+1)^{3/2}-n^{3/2}
       =\frac32x^{1/2}-\frac1{64}x^{-3/2}+O(x^{-7/2}).
\]
Consequently the combined summand in \eqref{et:critical-cubic} is
$-\tfrac12x^{-3}+O(x^{-5})$, proving absolute convergence.
In the exact cutoff formula,
$N^{3/2}T_N(3/2)^3\to8$, $T_N(3/2)^2\to0$, and
$H_N^{(3)}\to\zeta(3)$. These limits give the stated critical value.
\end{proof}


\section{Fifth and sixth powers: finite coordinate spaces and explicit identities}
\label{tp:sec:examples}

We now make the all-power theorem concrete.  Retain the classical
constants
\begin{align}
 Q_5&=4\zeta(2)\zeta(3)-10\zeta(5),\label{tp:eq:Q5}\\
 Q_7&=51\zeta(3)\zeta(4)+10\zeta(2)\zeta(5)
                       -\frac{301}{4}\zeta(7).\label{tp:eq:Q7}
\end{align}
The latter is $\sum_{n\geq0}A_n^4$, by Hoffman's quartic evaluation
with the $n=0$ term included \cite{HoffmanTails,RepoQuartic}.

\subsection{The lower-power inputs}

The universal recurrence starts with
\begin{align}
 H_1(t)&=-\frac{t^2}{4},\label{tp:eq:H1}\\
 H_2(t)&=3\zeta(3)t+2I^2(\cosh(t/2)G_1(t)),\label{tp:eq:H2}\\
 H_3(t)&=3I^2(\cosh(t/2)G_2(t))
      -\frac{17}{8}\zeta(4)t^2-\frac{\zeta(2)}{48}t^4
      +\frac{t^6}{2880},\label{tp:eq:H3}\\
 H_4(t)&=Q_7t+\frac{Q_5t^3}{6}
      +4I^2(\cosh(t/2)G_3(t))-I^4H_2(t).\label{tp:eq:H4}
\end{align}
These reproduce the existing quadratic, cubic, and quartic generators.
For example $\mathcal Z_{2,-1}=2\zeta(2,3)+\zeta(5)=-Q_5$ and
$\mathcal Z_{e,0}=\zeta(2)^e$ account for their exceptional terms.

\subsection{The quintic generator and its eight coordinates}

Define the one new weight-eight combination
\begin{equation}\label{tp:eq:K8}
 K_8=\frac56\mathcal Z_{3,-2}+\frac{289}{32}\zeta(8),
 \quad
 \mathcal Z_{3,-2}=
 6\zeta(2,2,4)+3\zeta(2,6)+3\zeta(4,4)+\zeta(8).
\end{equation}
This is a specified convergent MZV coordinate, not a claim of a new
transcendental constant or of its irreducibility.

\begin{theorem}[Every regular quintic moment]
\label{tp:thm:five}
For $|t|<2\pi$,
\begin{equation}\label{tp:eq:H5}
 \boxed{\begin{aligned}
 H_5(t)={}&5I^2(\cosh(t/2)G_4(t))
            -\frac53I^4H_3(t)-K_8t^2
            -\frac{19}{64}\zeta(6)t^4\\
          &+\frac{\zeta(4)}{4320}t^6
            +\frac{\zeta(2)}{241920}t^8
            -\frac{t^{10}}{43545600}.
 \end{aligned}}
\end{equation}
All the moments $M_{5,2r+1}$ lie in
\begin{equation}\label{tp:eq:space5}
 \operatorname{span}_{\mathbb Q}
 \{K_8,\zeta(6),Q_7,Q_5,\zeta(2),\zeta(3),\zeta(4),1\}.
\end{equation}
The coefficient matrix of the first eight moments in this ordered
list has determinant
\begin{equation}\label{tp:eq:rank5}
 \frac{672072587373661}{748313054183581286400}\ne0.
\end{equation}
Thus eight is the exact rank of these specified rational coefficient
rows.  No arithmetic independence of the eight evaluated constants
is asserted.
\end{theorem}

\begin{proof}
At $d=5$, the two Bernoulli coefficients in \eqref{tp:eq:generator}
are $5/3$ and $-1/6$.  Insert the anchors
$\mathcal Z_{3,-2}$, $\mathcal Z_{3,0}=\zeta(2)^3$, and
$\mathcal Z_{1,-2j}=\zeta(2j+2)$.  Euler's even-zeta evaluations
give $\zeta(2)^3=35\zeta(6)/8$ and
$\zeta(2)^4=175\zeta(8)/24$.  Finally insert
$I^8H_1=-t^{10}/7257600$.  Collecting these finite terms gives
\eqref{tp:eq:H5}.  The span follows from the lower-power generators.
Exact rational coefficient extraction gives \eqref{tp:eq:rank5};
the complete matrix and an executable determinant calculation accompany
the article.
\end{proof}

The first three formulas, useful independently of the generator, are
\begin{align}
 M_{5,1}={}&-K_8+\frac52Q_7,\label{tp:eq:M51}\\
 M_{5,3}={}&\frac14K_8+\frac5{12}Q_5-\frac5{12}Q_7
           +\frac{15}{2}\zeta(3)-\frac{85}{24}\zeta(4)
           -\frac{57}{32}\zeta(6),\label{tp:eq:M53}\\
 M_{5,5}={}&-\frac7{48}K_8-\frac5{72}Q_5+\frac{47}{288}Q_7
           -\frac1{12}\zeta(2)-\frac52\zeta(3)
           +\frac{157}{144}\zeta(4)\notag\\
          &+\frac{95}{64}\zeta(6)-\frac1{12}.
          \label{tp:eq:M55}
\end{align}
The first two are unsubtracted sums.  The third is
$\sum_{n\geq0}(x_n^5A_n^5-1)$.

\begin{corollary}[Two ordinary quintic evaluations]
\label{tp:cor:explicit}
With $x=x_n$ and $A=A_n$, both series below converge absolutely:
\begin{align}
 &\sum_{n\geq0}\bigl((23x+232x^3+240x^5)A^5-240\bigr)\notag\\
 &\qquad=80Q_5-20\zeta(2)+1140\zeta(3)-560\zeta(4)
                         -57\zeta(6)-20,\label{tp:eq:short5}\\[3pt]
 &\boxed{\begin{aligned}
 &\sum_{n\geq0}\left[
 (195x+4056x^3+11568x^5+5376x^7)A^5
                  -5376x^2-9328\right]\\
 &\hspace{12mm}=-748\zeta(2)+7884\zeta(3)-5380\zeta(4)
                  +171\zeta(6)-748.
 \end{aligned}}\label{tp:eq:main5}
\end{align}
\end{corollary}

\begin{proof}
The first identity is $23M_{5,1}+232M_{5,3}+240M_{5,5}$.
Its $K_8$ and $Q_7$ coefficients cancel.  The second is
$195M_{5,1}+4056M_{5,3}+11568M_{5,5}+5376M_{5,7}$.
The three coefficients of $K_8,Q_7,Q_5$ cancel.  Its subtraction
polynomial is $5376(x^2-5/12)+11568=5376x^2+9328$.

There is a particularly small independent MZV check of the second
formula.  The finite-cutoff recursion \eqref{tp:eq:independentT}
reduces its left side to
\begin{align*}
 &-748\zeta(2)+5256\zeta(2,1)-3840\zeta(2,2)
   +240\zeta(2,2,2)\\
 &\quad+120\bigl(\zeta(2,4)+\zeta(4,2)\bigr)
   +2628\zeta(3)-2500\zeta(4)+36\zeta(6)-748.
\end{align*}
Use $\zeta(2,1)=\zeta(3)$,
$\zeta(2,2)=3\zeta(4)/4$,
$\zeta(2,2,2)=3\zeta(6)/16$, and
$\zeta(2,4)+\zeta(4,2)=3\zeta(6)/4$.
These are Euler, all-twos, and elementary stuffle identities.
The displayed result follows.  Absolute convergence follows directly
from \eqref{tp:eq:centeredasy}.
\end{proof}

The values of \eqref{tp:eq:short5} and \eqref{tp:eq:main5} begin,
respectively,
\[
 456.5518108340607851715385608386\ldots,
 \qquad1849.672608741185304145251799308\ldots.
\]
These decimals are diagnostic illustrations, not inputs to the proof.

\subsection{The sextic generator and eleven coordinates}

Define
\begin{equation}\label{tp:eq:new6}
 R_{11}=\mathcal Z_{6,1},\qquad
 K_9=\mathcal Z_{2,-5}-5\mathcal Z_{4,-1},\qquad
 J_7=\mathcal Z_{2,-3}=2\zeta(2,5)+\zeta(7).
\end{equation}
Formula \eqref{tp:eq:anchor} specifies every term of these finite
convergent MZV combinations.

\begin{theorem}[Every regular sextic moment]
\label{tp:thm:six}
For $|t|<2\pi$,
\begin{equation}\label{tp:eq:H6}
 \boxed{\begin{aligned}
 H_6(t)={}&R_{11}t+6I^2(\cosh(t/2)G_5(t))
          -\frac52I^4H_4(t)+\frac12I^8H_2(t)\\
         &+\frac{K_9}{12}t^3+\frac{J_7}{240}t^5
          -\frac{Q_5}{10080}t^7.
 \end{aligned}}
\end{equation}
All $M_{6,2r}$ belong to the rational span of the ordered list
\begin{equation}\label{tp:eq:space6}
 (R_{11},K_9,J_7,K_8,\zeta(6),Q_7,Q_5,
                       \zeta(2),\zeta(3),\zeta(4),1).
\end{equation}
The first eleven coefficient rows have determinant
\begin{equation}\label{tp:eq:rank6}
 -\frac{557325343494359097538722307}
         {135927478735611867783198720000000}\ne0.
\end{equation}
This is again a formal coefficient-rank statement.
\end{theorem}

\begin{proof}
The nonzero coefficients at $d=6$ in \eqref{tp:eq:generator} are
$5/2$ and $-1/2$.  The $k=1$ anchor is $R_{11}$; the remaining
negative-index anchors combine to $K_9,J_7$, and
$\mathcal Z_{2,-1}=-Q_5$.  This gives \eqref{tp:eq:H6} directly.
The span and determinant follow by rational coefficient extraction.
\end{proof}

For example,
\begin{align}
 M_{6,0}&=R_{11},\notag\\
 M_{6,2}&=-2K_8+\frac16K_9+5Q_7-\frac1{12}R_{11},\notag\\
 M_{6,4}&=\frac1{10}J_7+\frac25K_8-\frac1{12}K_9
       +\frac12Q_5-\frac54Q_7+\frac7{240}R_{11}\notag\\
       &\quad+9\zeta(3)-\frac{17}{4}\zeta(4)
                       -\frac{171}{80}\zeta(6).
                       \label{tp:eq:M6examples}
\end{align}
One further elimination gives an ordinary sixth-power identity:
\begin{equation}\label{tp:eq:main6}
 \begin{split}
 &\sum_{n\geq0}\left[
 (35+3360x^2+29792x^4+50176x^6+16128x^8)A^6
                    -16128x^2-42112\right]\\
 &\quad=9184Q_5-64Q_7-2928\zeta(2)+111216\zeta(3)
          -55368\zeta(4)-9576\zeta(6)-2928.
 \end{split}
\end{equation}
It is seven times the combination with coefficients
$(5,480,4256,7168,2304)$ of $M_{6,0},\ldots,M_{6,8}$.
Its value is $33200.31747342681779918605867754\ldots$.

\section{The resonant parity and all Laurent coefficients}
\label{tp:sec:resonance}

The regular moments do not determine spectral derivatives at the
interlaced poles.  The same explicit subtractions nevertheless give
the complete Laurent family, including all Stieltjes terms.  This
addresses the normalization part of the quartic report's Research
question R5 and applies to every power.

Use the convention
\begin{equation}\label{tp:eq:Stieltjes}
 \zeta(1+\varepsilon,a)=\frac1\varepsilon+
   \sum_{q\geq0}\frac{(-1)^q\gamma_q(a)}{q!}\varepsilon^q.
\end{equation}
For $d\geq1$, $r\geq0$, put $s_0=1-d-2r$ and
\[
 F_d(s_0+\varepsilon)=\frac{h_{d,r}}\varepsilon+
                 \sum_{q\geq0}\Lambda_{d,r,q}\frac{\varepsilon^q}{q!}.
\]

\begin{theorem}[Exact logarithmic sums at every resonance]
\label{tp:thm:Laurent}
For every $q\geq0$,
\begin{equation}\label{tp:eq:Laurent}
 \boxed{\begin{aligned}
 \Lambda_{d,r,q}={}&(-1)^q\sum_{n\geq0}(\log x_n)^q
 \left[x_n^{d+2r-1}A_n^d-
       \sum_{j=0}^{r}h_{d,j}x_n^{2r-1-2j}\right]\\
 &+\sum_{j=0}^{r-1}h_{d,j}
        \zeta^{(q)}(1+2j-2r,1/2)
       +(-1)^qh_{d,r}\gamma_q(1/2).
 \end{aligned}}
\end{equation}
The series is ordinary and absolutely convergent.  The Taylor series
of $F_d(s_0+\varepsilon)-h_{d,r}/\varepsilon$ has exact radius $2$.
\end{theorem}

\begin{proof}
The square bracket is $O(x_n^{-3})$. On a compact subset $K$ of
$\Re\varepsilon>-2$, its product with $x_n^{-\varepsilon}$ has
$q$th derivative bounded by
\[
 O_{K,q}\!\left(x_n^{-3-\min_{\varepsilon\in K}\Re\varepsilon}
                         (1+|\log x_n|)^q\right).
\]
This bound is summable, so the remainder and all fixed derivatives
converge normally there. Add
\[
 \sum_{j=0}^{r-1}h_{d,j}\zeta(1+2j-2r+\varepsilon,1/2)
 +h_{d,r}\bigl(\zeta(1+\varepsilon,1/2)-1/\varepsilon\bigr).
\]
This is the required pole-subtracted meromorphic family on a common
half-plane, hence by continuation near zero.  Differentiating there
gives \eqref{tp:eq:Laurent}, including the sign and normalization of
the Stieltjes term.  There are no poles in $|\varepsilon|<2$, while
$\varepsilon=-2$ is an actual pole with nonzero residue $h_{d,r+1}$.
This proves the radius assertion.
\end{proof}

At the first resonance the formula is especially short:
\begin{equation}\label{tp:eq:firstresonance}
 \Lambda_{d,0,q}=(-1)^q\left[
 \sum_{n\geq0}(\log x_n)^q
           \left(x_n^{d-1}A_n^d-x_n^{-1}\right)
          +\gamma_q(1/2)\right].
\end{equation}
All half-shift Stieltjes constants can be converted to the ordinary
constants $\gamma_q=\gamma_q(1)$ by
\begin{equation}\label{tp:eq:halfStieltjes}
 \gamma_q(1/2)=
 \frac{2(-1)^q(\log2)^{q+1}}{q+1}
 +2\sum_{j=0}^q\binom qj(-\log2)^{q-j}\gamma_j-\gamma_q.
\end{equation}
This follows by expanding
$\zeta(s,1/2)=(2^s-1)\zeta(s)$ at $s=1$.

The theorem gives a locally analytic, explicitly anchored generator in
the spectral displacement $\varepsilon$.  It does not identify a
finite arithmetic span as $r$ or the logarithmic order $q$ varies;
that is a further reduction problem.  In particular one must not
differentiate the discrete coefficient formulas for $M_{d,m}$ as
though they were identities at arbitrary complex $m$.

% END INCLUDED FILE: sections/02_all_powers.tex


% BEGIN INCLUDED FILE: sections/03_moment_completion.tex
% Research contribution fragment for the parent article.
% Source baseline: ProveIt da2792fc9cc8ef19c6a314be6c65e1437fbdc66b.
% All labels use mt10:. This is not a standalone final article.

\section{Completion of the smaller-order-ten reconstruction theorem}
\label{mt10:section}

Two incoming reports leave complementary parts of the same problem open.
The \emph{Alternating Tails and Spectral Defects} report \cite{RepoAlternating} proves the
smaller-order-ten statement outside one residue class modulo $2048$;
the \emph{Quartic Spectral Reconstruction} report \cite{RepoQuartic} proves it for all
sufficiently large larger orders. Their thresholds leave exactly one
matrix unaccounted for. We prove that this matrix is nonsingular by a
short exact certificate and hence complete the entire smaller-order-ten
family. The finite polynomial identities and the sufficient threshold
used from those reports are replayed independently below. Their
underlying reductions remain attributed to those reports.

This completion concerns formal rational coordinate rank. It implies
identities among actual polygamma sums, but does not assert arithmetic
independence of evaluated zeta values. The unrestricted unequal-order
conjecture remains open.

\subsection{Moment coordinates and the precise theorem}

For integers $s\ge2$ put
\[
 A_s(n)=\zeta(s,n+1)=\zeta(s)-H_n^{(s)},\qquad x_n=n+\tfrac12.
\]
For an integer $b\ge6$, initially define
\[
 D_{10,2b}(w)=\sum_{n\ge0}A_{10}(n)A_{2b}(n)x_n^{-w},
 \qquad \operatorname{Re}w>-2b-7.
\]
The centered Hurwitz expansion contains only powers
$x_n^{-2b-8-2j}$, so continuation by subtraction is regular at every
nonpositive even integer. Write $D_r=D_{10,2b}(-2r)$. Define the
mutually inverse finite transforms
\begin{align}
 C_j&=\sum_{r=0}^j\binom{2j+1}{2r}2^{-2(j-r)}D_r,
 \label{mt10:C}\\
 D_r&=\sum_{j=0}^r\frac{\binom{2r}{2j}}{2j+1}
                 B_{2r-2j}(1/2)C_j.\label{mt10:D}
\end{align}
Finite summation by parts gives, for every $j\ge b$,
\[
 C_j=\sum_{k=0}^{b-1}(M^{\infty}_{5,b})_{j-b,k}Z_k,
 \qquad Z_k=\zeta(2k),\quad Z_0=-\tfrac12,
\]
where, for $t\ge0$ and $0\le k<b$,
\begin{equation}\label{mt10:M}
 (M^{\infty}_{5,b})_{tk}
 =\mu_{t+k-4}\left\{
 \frac{(2b+2t-9)!}{(2b-2k)!}
 +\mathbf1_{k<5}\frac{(2t+1)!}{(10-2k)!}\right\},
 \qquad
 \mu_j=\begin{cases}B_{2j}/(2j)!,&j\ge0,\\0,&j<0.\end{cases}
\end{equation}
Let $M_{5,b}$ be its first $b$ rows. Equivalently, each entry is
\[
 B_{2(t+k-4)}\left\{
 \frac{\binom{2t+2}{10-2k}}{2t+2}
 +\frac{\binom{2b+2t-8}{2b-2k}}{2b+2t-8}\right\},
\]
where a negative Bernoulli index gives zero, and a binomial summand
is omitted unless its lower index lies between two and its upper
index. This last convention is essential.

The formal coordinate space here has independent labels
$C_0,\ldots,C_{b-1},Z_0,\ldots,Z_{b-1}$. Its evaluation map assigns
the indicated complex numbers to those labels. In these coordinates
the first $2b$ transformed moments have matrix
$\operatorname{diag}(I_b,M_{5,b})$. Since the diagonal of
\eqref{mt10:D} is $1/(2r+1)$, the first $2b$ original moments have
determinant $\det M_{5,b}/(4b-1)!!$.

\begin{theorem}[Every larger even order at smaller order ten]
\label{mt10:main}
For every integer $b\ge6$, $M_{5,b}$ is nonsingular. Consequently
\[
 D_{10,2b}(0),D_{10,2b}(-2),\ldots,D_{10,2b}(-2(2b-1))
\]
have full formal coordinate rank $2b$, and every subsequent centered
even moment is an explicit rational linear combination of this prefix.
\end{theorem}

\subsection{The five-dimensional contraction and its normalization}

For $0\le i,\ell\le4$ set
\[
 n_i=b-5+i,\qquad N_{i\ell}=2(n_i-\ell)+1,
 \qquad R_\ell=2b-8+2\ell,
\]
\[
 F_{i\ell}=\frac1{(2\ell+1)!}
       \prod_{r=0}^{2\ell-1}(2n_i-r)(2n_i+3+r).
\]
An empty product is one. Define finite integer-polynomial coefficients
$d_{ih}$ by the Newton expansion
\begin{equation}\label{mt10:newton}
 (2n_i+2-y)(4n_i+2-y)^{\underline{2i}}
   =\sum_{h=0}^{2i+1}d_{ih}y^{\underline h}.
\end{equation}
Thus explicitly
\[
 d_{ih}=\frac1{h!}\sum_{j=0}^h(-1)^{h-j}\binom hj
           (2n_i+2-j)(4n_i+2-j)^{\underline{2i}}.
\]
For independent variables $U_1,\ldots,U_9$, let
\begin{equation}\label{mt10:A}
 \mathcal A_{i\ell}=2F_{i\ell}+
 \frac{R_\ell!}{(2n_i+1)(2n_i+2)(2\ell+2)!N_{i\ell}!}
 \sum_{k=0}^i d_{i,2k+1}N_{i\ell}^{\underline{2k+1}}
                      U_{5-i+\ell+k}.
\end{equation}
Ratios of factorials have fixed index difference and mean the
corresponding finite rational products in $b$.

For $b\ge10$, the exact reduction is
\begin{equation}\label{mt10:ratio}
 \rho_{5,b}:=\frac{\det M_{5,b}}{\det R\det E\det H_{5,b}}
  =\frac{\det\mathcal A(b;2B_{2b-2},\ldots,2B_{2b-18})}
         {\det F(b)},
\end{equation}
where
\[
 R_{tt}=(2b+2t-9)!,\qquad E_{kk}=\frac1{(2b-2k)!},
 \qquad (H_{5,b})_{tk}=\mu_{t+k-4}.
\]
All three determinants in the denominator are nonzero.

Here are details sufficient to check the contraction without an
assumed inverse-Hankel formula. The monic polynomial rows
\[
 \pi_n(x)=\sum_{t=0}^n
 \frac{2(2n+2t+1)!(2n+1)!}
 {(4n+2)!(2n-2t+1)!(2t)!}x^t
\]
are orthogonal for the functional $\mathcal N(x^r)=\mu_{r+1}$:
\[
 \mathcal N(x^j\pi_n)=0\quad(0\le j<n),\qquad
 \mathcal N(x^n\pi_n)=h_n,
 \quad h_n=\frac{((2n+1)!)^2}{(4n+2)!(4n+3)!}>0.
\]
These follow by coefficient comparison in the elementary Pad\'e identity
\[
 e^tP_N(-t)-P_N(t)
 =\frac{(-1)^N(N!)^2}{(2N)!(2N+1)!}t^{2N+1}
        +O(t^{2N+2}),
\]
\[
 P_N(t)=\sum_{k=0}^N
       \frac{(2N-k)!N!}{(2N)!k!(N-k)!}t^k,
 \qquad N=2n+1.
\]
The coefficient cancellation through degree $2N$ is the finite binomial
identity $\sum_k(-1)^k\binom Nk(1+x)^{2N-k}=x^N(1+x)^N$;
the next coefficient is an elementary beta integral.

Replacing every row of $H_{5,b}+U^{\mathrm{pert}}$ by these monic
rows is a unit lower triangular operation. Its final five rows vanish
in columns $k\ge5$. The unperturbed retained entries, after column
reversal and division by
$c_n=(2n+1)(2n+2)!/(4n+2)!$, are precisely $F_{i\ell}$.
For the perturbation apply
\[
 \sum_{j=0}^N\binom Nj B_j(N-j)^{\underline h}
                  =N^{\underline h}B_{N-h}(1)
\]
to \eqref{mt10:newton}. Because $b\ge10$, every surviving index
$N-h$ is at least two. The even $h$ give vanishing odd Bernoulli
numbers; restoring the omitted term $j=1$, with $B_1=-1/2$, contributes
a second $F_{i\ell}$. This proves \eqref{mt10:A}, including both
factors of two. Carrying the same row operation out on $H_{5,b}$
alone gives \eqref{mt10:ratio}.

For completeness, $H_{5,b}$ is nonsingular: after conjugating by
$\operatorname{diag}(1,-1,1,\ldots)$ and splitting off its first
three indices, it has the form
\[
 \begin{pmatrix}ee^{\mathsf T}&B\\B^{\mathsf T}&-P\end{pmatrix},
 \qquad e=(0,0,1)^{\mathsf T},
\]
where $P$ is positive definite. Its entries are moments on the infinite
positive support $\{(2\pi m)^{-2}:m\ge1\}$. The first two rows
of $B$ have a nonsingular antitriangular two-by-two minor. Therefore
$ee^{\mathsf T}+BP^{-1}B^{\mathsf T}$ is positive definite, which
proves nonsingularity by the Schur complement. This argument also
shows that the unperturbed contraction $F$ is nonsingular.

\subsection{The inherited congruence, with its finite certificate replayed}

Define
\begin{align*}
 D_{10}(b)={}&4096b^{12}-110592b^{11}+1236992b^{10}-7303680b^9\\
 &+23705088b^8-38140416b^7+14925056b^6+24490080b^5\\
 &-4469984b^4-15838992b^3-2571948b^2+4244400b+1488375.
\end{align*}
The exact five-by-five polynomial identities are
\begin{align}
 \det F={}&\frac{16384}{4465125}
 (b-3)(b-2)^2(b-1)(2b-7)(2b-5)^2(2b-3)\notag\\
 &\qquad\cdot(4b-15)(4b-13)(4b-11)^2(4b-9)^2(4b-7)^2
 (4b-5)(4b-3)D_{10}(b),\label{mt10:Fdet}\\
 \frac{\det\mathcal A}{\det F}
   &=\frac{P_{10}(b;U_1,\ldots,U_9)}{
        138240(2b-7)(2b-5)D_{10}(b)},\label{mt10:P}
\end{align}
where $P_{10}$ is an integer polynomial. Equations \eqref{mt10:A}
and \eqref{mt10:P} define it exactly; its supplied sparse certificate
contains $10{,}403$ monomials, grouped into $439$ monomials in the
$U$ variables. Their coefficients have degree at most $24$ in $b$.

The replay forms the row denominators of \eqref{mt10:A}, constructs
the $120$ signed permutation products of the cleared matrix, and
cross-multiplies \eqref{mt10:P} in
$\mathbb Q[b,U_1,\ldots,U_9]$. Equality is checked coefficient by
coefficient. Thus polynomial integrality and the displayed determinant
identity are finite rational certificates, without interpolation.

Put
\[
 \mathcal B(x)=1+20068x-26384x^2+2816x^3-19712x^4
                         +4096x^5-24576x^6.
\]
The uniform Bernoulli congruence is
\begin{equation}\label{mt10:B16}
           2B_{2m}\equiv\mathcal B(m)\pmod {65536}
                  \quad(m\ge12).
\end{equation}
All congruences are in $\mathbb Z_{(2)}$, the rationals with odd
denominator. One direct proof starts with Faulhaber's formula at
$256$. The terms other than $256B_{2m}$ have valuation at least
$23$; after division by $128$, the even summands vanish at the
required precision because $2m\ge24$. For odd $u$, $8\mid u^2-1$,
so expansion of $u^{2m}=(1+(u^2-1))^m$ gives
\[
 2B_{2m}\equiv\frac1{128}
       \sum_{u=1,3,\ldots,255}\sum_{j=0}^7
                    \binom mj(u^2-1)^j\pmod {65536}.
\]
The omitted terms have valuation at least $3\cdot8-7=17$.
Expanding this finite average in $\mathbb Q[m]$ and reducing every
coefficient modulo $65536$ gives \eqref{mt10:B16}. The replay
checks this coefficientwise, including the denominators of
$\binom mj$. The bound $v_2(B_{2j})\ge-1$ used in the Faulhaber
step follows inductively from the odd-denominator recurrence
\eqref{mt10:Bernoulli-recurrence} below.

The finite polynomial substitution then gives
\begin{equation}\label{mt10:congruence}
 P_{10}\bigl(b;\mathcal B(b-1),\ldots,\mathcal B(b-9)\bigr)
                  \equiv32\mathcal R(b)\pmod {65536},
\end{equation}
where
\[
 \mathcal R(b)=-512b^8+128b^7-928b^6-816b^5+464b^4
                         +352b^3+118b^2+217b+694.
\]
The replay verifies every coefficient of \eqref{mt10:congruence}
using modular polynomial multiplication. The linear coefficient of
$\mathcal R$ is odd and every coefficient of degree at least two
is even. Consequently
\[
 v_2(\mathcal R(x)-\mathcal R(y))=v_2(x-y)\quad(x\ne y).
\]
Indeed the quotient by $x-y$ is odd. Direct substitution gives
$\mathcal R(962)\equiv0\pmod {2048}$, so $962$ is its unique
root modulo $2048$.

Since $D_{10}(b)$ is odd at every integer $b$, the denominator in
\eqref{mt10:P} has valuation $10$. For $b\ge21$ all nine Bernoulli
indices are in the range of \eqref{mt10:B16}. If
$b\not\equiv962\pmod {2048}$, equations \eqref{mt10:B16}--
\eqref{mt10:congruence} therefore prove
\begin{equation}\label{mt10:valuation}
                 v_2(\rho_{5,b})=v_2(b-962)-5.
\end{equation}
For $6\le b\le20$, direct rational determinants of the original
matrix \eqref{mt10:M} give the respective valuations
\[
 -4,-4,-6,-5,-2,-5,-4,-5,-3,-5,-4,-5,-1,-5,-4.
\]
The direct matrices, their determinants, and their normalizations are
included in the replay record. In particular, the contraction is never
used below its valid threshold $b\ge10$.

\subsection{A transparent sufficient large-order threshold}

We specialize the incoming eventual-rank proof to $a=5$ and replace
its large rational constant by an elementary upper bound. This makes
the threshold independently reviewable.

Define
\[
 J(b)=\prod_{i=0}^4(2n_i+1)(2n_i+2),\qquad
 E(b)=\prod_{i,\ell=0}^4
             \prod_{r=1}^{\max(0,2i-4\ell-1)}(R_\ell+r),
 \quad D(b)=10!J(b)E(b).
\]
Here $\deg J=10$, $\deg E=20$, and $D\mathcal A$ has integer
polynomial entries. Put $\mathscr P=\det(D\mathcal A)$ and let
$\mathscr P_{[5]}$ be its homogeneous part of degree five in the
$U$ variables. Every entry has degree at most $46$ in $b$; hence
$\deg_b\mathscr P\le230$. The largest coefficient norm among the
$25$ entries, with all coefficients counted in
$\mathbb Z[b,U_1,\ldots,U_9]$, is
\[
 E^*=602847821074996042728613866361551228764160<10^{42}.
\]
The separate verifier expands the $25$ entries, records all $25$
norms and degrees, and reproduces this integer. The full determinant
coefficient norm is at most $K=5!(E^*)^5<10^{213}$.

Let $F_0=(2b-18)!$ and $V_j=U_jF_0/(2b-2j)!$. From
\eqref{mt10:newton},
$d_{i1}=-(i+1)(4n_i+2)^{\underline{2i}}$. Cancelling factorials
in \eqref{mt10:A} gives
\[
 \mathcal A-2F=
 \operatorname{diag}\left(\frac1{(2n_i+1)(2n_i+2)}\right)
 L(V_{5-i+\ell})_{0\le i,\ell\le4}
 \operatorname{diag}\left(\frac{R_\ell!}{F_0(2\ell+2)!}\right),
\]
where $L_{i,i-k}=d_{i,2k+1}$ is lower triangular. Therefore
\[
 \mathscr P_{[5]}=-\Gamma(b)\det(V_{i+j+1})_{0\le i,j\le4},
\]
with the explicit positive factor
\begin{align*}
 \Gamma(b)={}&\kappa J(b)^4E(b)^5
 \prod_{i=0}^4(4n_i+2)^{\underline{2i}}
 \prod_{\ell=0}^4\prod_{r=1}^{10+2\ell}(2b-18+r),\\
 \kappa={}&\frac{5!(10!)^5}{2!4!6!8!10!}
          =14932726824960000000>10^{19}.
\end{align*}
Every linear factor is at least $b$ for $b\ge20$, so
$\Gamma(b)\ge\kappa b^{230}$.

At the actual doubled Bernoulli values, Euler's even-zeta formula
and row/column reversal yield
\[
 |\mathscr P_{[5]}|
 =\frac{\Gamma(b)(4F_0)^5}{(2\pi)^{10(b-5)}}
       \det\bigl(\zeta(2b-18+2i+2j)\bigr)_{0\le i,j\le4}.
\]
The determinant on the right is a positive moment determinant on
the points $m^{-2}$. Keeping the support points $1,1/4,1/9,1/16,1/25$
in Cauchy--Binet proves
\begin{gather*}
 \det\bigl(\zeta(2b-18+2i+2j)\bigr)\ge V\,120^{-2b+18},\\
 V=\prod_{1\le r<s\le5}(r^{-2}-s^{-2})^2
        =\frac{2401}{36864000000000000}>10^{-14}.
\end{gather*}

For $b\ge30$, all nine actual variables have absolute value at most
\[
       T_b=\frac{8(2b)^{16}F_0}{(2\pi)^{2b-18}}\ge1.
\]
Thus $|\mathscr P_{[\le4]}|\le Kb^{230}T_b^4$.
After cancellation of the common powers of $2\pi$, using
$1<2\pi<8$ gives
\begin{equation}\label{mt10:dominance}
 \frac{|\mathscr P_{[\le4]}|}{|\mathscr P_{[5]}|}
 \le C\frac{b^{64}960^{2b}}{(2b-18)!},\qquad
 C=\frac{K2^{142}}{\kappa4^5V120^{18}}<10^{251}<10^{256}.
\end{equation}
The powers here can be checked directly: the uncancelled power of
$2\pi$ is $2b+22$, and $(2b)^{64}$ comes from $T_b^4$.
The bounds $K<10^{213}$, $2^{142}<10^{43}$, $\kappa>10^{19}$,
and $V>10^{-14}$ already give the stated coarse bound on $C$.

The remaining comparison is an exact integer inequality:
\begin{equation}\label{mt10:base}
       2^{384}\,10^{256}\,1920^{64}\,960^{3840}<3822!.
\end{equation}
It is verified by integer multiplication, with both sides explicitly
specified. Moreover, for $b\ge1920$,
\[
 \frac{(b+1)^{64}960^{2b+2}/(2b-16)!}
      {b^{64}960^{2b}/(2b-18)!}
 =\left(1+\frac1b\right)^{64}
            \frac{960^2}{(2b-16)(2b-17)}<\frac34.
\]
Indeed $(1+1/b)^{64}<3$ and both denominator factors are at least
$b\ge1920=2\cdot960$. It follows from
\eqref{mt10:dominance}--\eqref{mt10:base} that the leading homogeneous
term cannot be cancelled when $b\ge1920$. Since $D(b)>0$ there,
\eqref{mt10:ratio} proves nonsingularity throughout that range.

\subsection{The single missing matrix has a small exact certificate}

The congruence covers every $6\le b<1920$ except $b=962$.
At this value define
\[
 S=\operatorname{diag}(2,1,64,1/8,2)\,
 \mathcal A\,\operatorname{diag}(1,1,1/32,1/16,1/512).
\]
Thus $\det S=2^{-13}\det\mathcal A$.
Every coefficient of this affine matrix in $U_1,\ldots,U_9$ has
$2$-adic valuation at least $-3$. Therefore its reduction modulo
$2^{20}$ is determined by the inputs modulo $2^{23}$.

The nine input residues, in increasing order of the index $j$, are
\[
 (660565,5568257,602509,3180025,6422085,3501169,
                         2681469,1596009,6743605)\pmod {2^{23}}.
\]
They are obtained by the elementary Bernoulli recurrence
\begin{equation}\label{mt10:Bernoulli-recurrence}
 (2m+1)(2B_{2m})
   =2m-1-\sum_{j=1}^{m-1}\binom{2m+1}{2j}(2B_{2j}),\qquad m\ge1.
\end{equation}
Every divisor $2m+1$ is odd. The recurrence is valid modulo every
power of two and produces the residues through $m=961$ without
floating arithmetic. A separate exact rational Bernoulli evaluation
agrees with all nine inputs.

Substitution in \eqref{mt10:A} gives the following matrix modulo
$2^{20}=1048576$:
\[
 S\equiv
 \begin{pmatrix}
 917951&442211&476112&551248&968980\\
 679904&45903&922072&654544&644616\\
 797935&540216&1028224&572168&196164\\
 491016&385277&695495&798931&823817\\
 917289&206367&412992&727792&439172
 \end{pmatrix}\pmod {1048576}.
\]
The determinant of this displayed integer matrix is
$524288=2^{19}$ modulo $2^{20}$. Hence
\[
 v_2(\det S)=19,\qquad
 v_2(\det\mathcal A)=32,\qquad
 v_2(\det F)=26,\qquad v_2(\rho_{5,962})=6.
\]
The value $v_2(\det F)=26$ follows immediately from
\eqref{mt10:Fdet}: its only even factors beyond $2^{14}$ are
$(b-2)^2=960^2$, of valuation $12$.
An independent check evaluates $\det\mathcal A/\det F$ directly
as a rational number; it gives the same valuation. The original
integer numerator also satisfies
$P_{10}(962;2B_{1922},\ldots,2B_{1906})\equiv65536\pmod {131072}$.

This is a finite nonvanishing certificate, not a numerical
approximation. It closes the only missing index below the sufficient
threshold. Together with the preceding congruence and dominance
arguments, it proves Theorem~\ref{mt10:main}.

\subsection{Ordinary convergent identities supplied by the completed family}

For $b\ge6$ and $0\le k<b$, let
$\lambda^{(b,k)}_t=(M_{5,b}^{-1})_{kt}$, $0\le t<b$, and define
the finite odd polynomial and its centered difference
\[
 K_{b,k}(u)=\sum_{t=0}^{b-1}\lambda^{(b,k)}_tu^{2b+2t+1},
 \qquad Q_{b,k}(x)=K_{b,k}(x+\tfrac12)-K_{b,k}(x-\tfrac12).
\]
The centered difference is precisely the coefficient transform
\eqref{mt10:C}, so its continued moment sum equals $Z_k$.
For the ordinary formulation put
\[
 \eta_{s,0}=\frac1{s-1},\qquad
 \eta_{s,j}=\frac{(s)_{2j-1}B_{2j}(1/2)}{(2j)!}\quad(j\ge1),
 \qquad c_{b,\ell}=\sum_{j=0}^{\ell}\eta_{10,j}\eta_{2b,\ell-j}.
\]
Let $T_{b,k}(x)$ be the polynomial consisting of the nonnegative
powers in
\[
             Q_{b,k}(x)x^{-2b-8}\sum_{\ell\ge0}c_{b,\ell}x^{-2\ell}.
\]
Only finitely many coefficients enter it, and its degree is at most
$2b-10$. Then the completed family gives the exact identities
\begin{equation}\label{mt10:ordinary}
 \boxed{\displaystyle
 \sum_{n=0}^{\infty}\left\{
 Q_{b,k}(x_n)
  \frac{\psi^{(9)}(n+1)\psi^{(2b-1)}(n+1)}{9!\,(2b-1)!}
       -T_{b,k}(x_n)\right\}=\zeta(2k).}
\end{equation}
The summands are $O(x_n^{-2})$. Every removed nonnegative even
power has centered continued sum $\zeta(-2j,1/2)=0$, so the
ordinary absolutely convergent sum has exactly the same value as
the moment combination. Formula \eqref{mt10:ordinary} is a finite
rational algorithm at each $b,k$, with nonsingularity now proved
for the entire parameter range. At $k=0$ its value is the rational
number $-1/2$.

\paragraph{Status and further research.}
The new logical step is the exact $b=962$ certificate and the
combination of the two independently audited ranges. The source
congruence, the finite contraction, and the eventual dominance
mechanism are retained with attribution. The same strategy is
available for smaller order twelve, but its finite cases and modular
determinant have not been proved here. A uniform argument for all
unequal even orders, or even the full gap-six family, is still a
substantive open problem. Higher-minor positivity cannot simply be
inferred from the known pairwise ratio inequalities; the determinant
sign already changes in some families.

\paragraph{Reproducible artifacts.}
\texttt{completion\_certificate.py} proves the missing case from
the displayed five-square matrix using no external project files.
\texttt{verify\_eventual\_ten.py} rederives all $25$ entry norms and
the exact comparison \eqref{mt10:base}.
\texttt{verify\_p10\_dependencies.py}, its attributed algebra helper,
and \texttt{p10\_dependency\_certificate.json} replay the full inherited
polynomial identity, modular substitution, and all small base matrices.
The three output records distinguish full replay from the deliberately
optional quick mode. None uses floating-point evidence for a theorem.

% END INCLUDED FILE: sections/03_moment_completion.tex


% BEGIN INCLUDED FILE: sections/04_harmonic_gaps.tex
% Fragment for the parent article. Requires amsmath, amssymb, amsthm.
% No documentclass or standalone bibliography is included.

\section{Several transverse harmonic marks and Gamma--Stieltjes gaps}
\label{gm:section}

The harmonic-gap continuation \cite{gm:gaps-source} leaves a precise
question after its single-marker theorem: can two independent transverse
derivatives be taken at adjacent deleted indices, while retaining the
prescribed relative continuation at nonintegral endpoints?  We answer its
first target,
\(\left.\partial_u\partial_v\mathcal H_{a,b}(s,u,v,t)\right|_{u=v=0}\),
and every analogous cluster of identical marks.  A finite set-partition
formula also determines the fully symmetric part of arbitrarily different
marks.  The partition coefficients are the classical symmetric-sum
coefficients of Hoffman \cite{gm:Hoffman1992,gm:Hoffman2018}.  The point
requiring proof here is their local realization by decorated relative
quotients, including products of Hurwitz primitives and their simultaneous
resonant derivatives.  An individual ordered word with unequal marks
retains information that the symmetric formula does not determine.

\subsection{Normalization and closure under finite decorations}

Throughout this section, \(\Re a,\Re b>0\), and all powers use the
analytic principal logarithm on the right half-plane.  For sufficiently
decaying lattice functions, define
\[
 Z_c[f_1,\ldots,f_d]
 =\sum_{n_1>\cdots>n_d\ge0}\prod_{j=1}^df_j(n_j+c),
 \qquad Z_c[\varnothing]=1.
\]
The relative quotient is fixed by
\begin{equation}\label{gm:relative}
 \mathcal H_{a,b}[w]
 =Z_b[w]-\sum_{j=1}^{|w|}Z_a[w_{\le j}]
                                  \mathcal H_{a,b}[w_{>j}],
 \qquad \mathcal H_{a,b}[\varnothing]=1.
\end{equation}
For power letters \(f_s(x)=x^{-s}\), write
\(\mathcal H_{a,b}(s_1,\ldots,s_d)\).  This is the relative
interpolant already used in the source; no new endpoint interpolation is
being chosen.

For clarity, its joint spectral entireness follows from the published
Ihara--Nakamura--Yamamoto theory \cite[Corollary~4.2 and
Proposition~5.3]{gm:INY}.  With deconcatenation convolution denoted by
\(*\), put \(\Phi_c(w)=\Psi(w;c-1)\) in their notation.  Their
formula is \(\Phi_c=Z_c^{-1}*Z_1\); consequently
\[
 \mathcal H_{a,b}=Z_a^{-1}*Z_b=\Phi_a*(\Phi_b)^{-1}.
\]
On each fixed word, convolution and inversion are finite polynomial
operations in the entire \(\Phi_c\)-values.  The harmonic product law
also follows from the character property
\cite[Theorem~5.2]{gm:INY}.  These identities initially hold on an open
absolute-convergence domain and then continue analytically.

Use the completed Hurwitz function
\begin{equation}\label{gm:Q}
 Q(u,x)=\zeta(u,x)-\frac1{u-1},\qquad
 Q(u,x+1)=Q(u,x)-x^{-u}.
\end{equation}
It is entire in \(u\) and holomorphic for \(\Re x>0\).
Write \(Q^{[r]}\) and \(\zeta^{[r]}\) for derivatives in the first
argument.  The Stieltjes convention is
\begin{equation}\label{gm:stieltjes}
 Q(1+\varepsilon,x)
  =\sum_{r\ge0}\frac{(-1)^r\gamma_r(x)}{r!}\varepsilon^r,
 \qquad Q^{[r]}(1,x)=(-1)^r\gamma_r(x).
\end{equation}

A decorated letter is \(x^{-s}\) multiplied by a polynomial in
finitely many \(Q(u,x)\) and \(Q(u,x+1)\).  Its \(Z_c\) and
\(\mathcal H_{a,b}\) are first defined by their actual convergent
series and \eqref{gm:relative}.  A common right half-plane in all
spectral variables exists.

\begin{lemma}[One tail with arbitrary other decorations]
\label{gm:one-tail}
Let \(w=(f_1,\ldots,f_d)\).  Denote by \(M_i(u)w\) the word with
\(f_i(x)\) replaced by \(f_i(x)Q(u,x)\), by \(I_j(u)w\) the
insertion of \(f_u\) after position \(j\), and by \(C_j(u)w\) the
replacement of \(f_j\) by \(f_jf_u\).  Then
\begin{equation}\label{gm:insertion}
 \mathcal H_{a,b}[M_i(u)w]
 =Q(u,a)\mathcal H_{a,b}[w]
  +\sum_{j=0}^{i-1}\mathcal H_{a,b}[I_j(u)w]
  +\sum_{j=1}^{i}\mathcal H_{a,b}[C_j(u)w].
\end{equation}
The other letters may already carry arbitrary finite decorations.
\end{lemma}

\begin{proof}
First replace \(Q\) by \(\zeta\) and work where all relevant sums
converge absolutely.  In the infinite decorated sum, the additional
Hurwitz-tail index is at least the index in position \(i\).  Its strict
positions are the \(i\) slots preceding that position; its equality
cases are the first \(i\) positions.  This gives the insertion and
merger terms with \(Z_c\) in place of \(\mathcal H_{a,b}\), without
an endpoint term.

Expand the relative recursion and group terms by their cut in the
original word.  At a cut before position \(i\), the insertions and
mergers contained in the prefix sum to
\(\zeta(u,a)Z_a[\text{prefix}]\), by the ordinary singleton stuffle
identity.  At a cut at or after position \(i\), the same counting
identity gives the decorated prefix.  The remaining suffix terms are
exactly the recursion for the right side of \eqref{gm:insertion}.
Triangular uniqueness in the word length proves that formula with
\(\zeta\).  Subtracting \(\mathcal H_{a,b}[w]/(u-1)\) gives the
stated completed formula.  This argument only uses pointwise products
of the other letters, so it applies unchanged when they are decorated.
\end{proof}

\begin{theorem}[Entire continuation of every finite decoration]
\label{gm:closure}
Every decorated relative word above continues jointly entirely in all
its spectral variables.  It has a terminating finite expansion into
bare relative words at affine sums of spectral arguments, with
coefficients polynomial in entire endpoint functions \(Q(U,a)\).
All fixed independent spectral derivatives are legitimate.
\end{theorem}

\begin{proof}
Use \eqref{gm:Q} to remove shifted factors.  Apply
Lemma~\ref{gm:one-tail} to one unshifted factor.  Every resulting
relative word contains one fewer \(Q\) factor: an insertion or merger
adds a bare power, while the remaining decorations are unchanged.
Induction ends with the entire bare relative words.  Each step agrees
with the initial convergent quotient, so the resulting entire expression
is its analytic continuation.  This proves both existence and
independence of the elimination order.  In the initial domain, Hurwitz
asymptotics uniformly on a slightly larger spectral compact set justify
termwise differentiation; elsewhere the derivatives follow from the
proved entire identities.  In particular the argument does not infer an
analytic identity from equality at integer endpoint differences.
\end{proof}

\subsection{The local symmetric-sum theorem}

For a finite separable polynomial \(P(y,z)\), define linearly
\begin{equation}\label{gm:E}
 \mathcal E_{a,b}^{s,t}[A(y)B(z)]
       =\mathcal H_{a,b}[f_sA,f_tB].
\end{equation}
We use this map for polynomials in the completed gap primitives and
their derivatives.  Theorem~\ref{gm:closure} specifies its continuation.
It also has an explicit ordinary-series meaning.  With
\(F(y,z)=y^{-s}z^{-t}P(y,z)\),
\begin{align}\label{gm:ordinary-E}
 \mathcal E_{a,b}^{s,t}[P]
 ={}&\sum_{n>j\ge0}F(n+b,j+b)
     +\sum_{0\le n\le j}F(n+a,j+a)\notag\\
    &-\sum_{n,j\ge0}F(n+a,j+b),
\end{align}
whenever these sums converge absolutely.  This follows by expanding
the depth-two case of \eqref{gm:relative}; the two \(a,a\) terms
combine to the weak opposite ordering.

Put
\begin{equation}\label{gm:gap}
 g_u(y,z)=Q(u,z+1)-Q(u,y).
\end{equation}
Taking the difference between the insertions at the two boundary
letters in Lemma~\ref{gm:one-tail}, and using \eqref{gm:Q}, gives
\begin{equation}\label{gm:local-insertion}
 \mathcal E_{a,b}^{s,t}[g_u]=\mathcal H_{a,b}(s,u,t).
\end{equation}
The same assertion holds with arbitrary additional boundary
decorations.

\begin{lemma}[Products of gap primitives]
\label{gm:ordered-partitions}
For independent complex \(u_1,\ldots,u_k\), let
\(U_B=\sum_{i\in B}u_i\).  Then
\begin{equation}\label{gm:ordered-partition-formula}
 \mathcal E_{a,b}^{s,t}\!\left[\prod_{i=1}^kg_{u_i}\right]
 =\sum_{(B_1,\ldots,B_\ell)}
      \mathcal H_{a,b}(s,U_{B_1},\ldots,U_{B_\ell},t),
\end{equation}
where the sum ranges over all ordered set partitions of
\(\{1,\ldots,k\}\).  The identity is jointly entire.
\end{lemma}

\begin{proof}
After some marked positions have been inserted, write their variables
as \(x_1,\ldots,x_\ell\), with \(x_0=y\) and
\(x_{\ell+1}=z\).  The exact telescoping identity
\begin{equation}\label{gm:split}
 g_v(y,z)=\sum_{h=0}^\ell g_v(x_h,x_{h+1})
                              +\sum_{h=1}^\ell x_h^{-v}
\end{equation}
follows from \eqref{gm:Q}.  It is an identity of primitive functions,
without an integer-lattice assumption on the external endpoints.
By \eqref{gm:local-insertion}, the first sum inserts the next label
in each possible strict slot, while the second merges it with each
existing block.  Starting with one marker, this constructs every
ordered set partition once.  All operations are identities between the
initial convergent quotients; Theorem~\ref{gm:closure} supplies their
joint continuation.
\end{proof}

For a set partition \(\pi\) of \(\{1,\ldots,k\}\), set
\[
 \mu(\pi)=\prod_{B\in\pi}(-1)^{|B|-1}(|B|-1)!.
\]

\begin{theorem}[Local symmetric-sum identity]
\label{gm:symmetric}
For arbitrary complex spectral arguments,
\begin{align}\label{gm:symmetric-formula}
 &\sum_{\sigma\in S_k}
       \mathcal H_{a,b}(s,u_{\sigma(1)},\ldots,u_{\sigma(k)},t)
 \notag\\
 &\qquad=\mathcal E_{a,b}^{s,t}\!\left[
        \sum_{\pi}\mu(\pi)\prod_{B\in\pi}g_{U_B}(y,z)\right].
\end{align}
\end{theorem}

\begin{proof}
Apply M\"obius inversion on the lattice of equality partitions to
\eqref{gm:ordered-partition-formula}.  A block \(B\) has M\"obius
coefficient \((-1)^{|B|-1}(|B|-1)!\).  This is the classical
symmetric-sum argument \cite{gm:Hoffman1992,gm:Hoffman2018};
Lemma~\ref{gm:ordered-partitions} proves that it acts on the present
local analytic expressions.  For example, at \(k=2\),
\[
 \mathcal H(s,u,v,t)+\mathcal H(s,v,u,t)
          =\mathcal E[g_ug_v-g_{u+v}].
\]
Every function used in the inversion has already been continued
jointly entirely.
\end{proof}

\subsection{Independent derivatives and identical marked clusters}

For integers \(r_i\ge0\) and centers \(\lambda_i\in\mathbb C\),
write
\[
 r_B=\sum_{i\in B}r_i,\qquad
 \lambda_B=\sum_{i\in B}\lambda_i,\qquad
 p_B(y,z)=Q^{[r_B]}(\lambda_B,z+1)
                              -Q^{[r_B]}(\lambda_B,y).
\]

\begin{corollary}[All independently marked symmetric sums]
\label{gm:marked}
The exact identity
\begin{align}\label{gm:marked-formula}
 &\left.\prod_{i=1}^k\partial_{u_i}^{r_i}
    \sum_{\sigma\in S_k}
      \mathcal H_{a,b}(s,u_{\sigma(1)},\ldots,u_{\sigma(k)},t)
                       \right|_{u_i=\lambda_i}\notag\\
 &\qquad=\mathcal E_{a,b}^{s,t}\!\left[
               \sum_\pi\mu(\pi)\prod_{B\in\pi}p_B(y,z)\right]
\end{align}
uses at most \(2^k-1\) different gap primitives.  At every merged
resonance \(\lambda_B=1\), the corresponding primitive is
\begin{equation}\label{gm:merged-stieltjes}
 p_B(y,z)=(-1)^{r_B}
                  \bigl(\gamma_{r_B}(z+1)-\gamma_{r_B}(y)\bigr).
\end{equation}
\end{corollary}

\begin{proof}
Differentiate the entire identity \eqref{gm:symmetric-formula}.
Within a block, every variable enters \(U_B\) with coefficient one,
so its derivative order is exactly \(r_B\), without an extra
multinomial coefficient.  Equation~\eqref{gm:stieltjes} proves
\eqref{gm:merged-stieltjes}.  Simultaneous resonances therefore require
no directional limit or separately chosen finite part.
\end{proof}

When all pairs \((\lambda_i,r_i)\) equal \((\lambda,r)\), each
permutation term is identical.  Put
\begin{align}\label{gm:Bell}
 p_j&=Q^{[jr]}(j\lambda,z+1)-Q^{[jr]}(j\lambda,y),\notag\\
 e_k(p_1,\ldots,p_k)
   &=[v^k]\exp\left(\sum_{j=1}^k\frac{(-1)^{j-1}}j p_jv^j\right).
\end{align}
The set-partition formula becomes the individual ordered identity
\begin{equation}\label{gm:identical}
 \left.\partial_{u_1}^r\cdots\partial_{u_k}^r
       \mathcal H_{a,b}(s,u_1,\ldots,u_k,t)
                         \right|_{u_1=\cdots=u_k=\lambda}
       =\mathcal E_{a,b}^{s,t}[e_k(p_1,\ldots,p_k)].
\end{equation}
These are independent derivatives taken before specialization.  A
total derivative in one common variable would give a different sum.
The coefficient in \eqref{gm:Bell} is a finite polynomial operation;
no convergence of an all-\(k\) marked generating series is asserted.
The first useful polynomials are
\begin{align}\label{gm:Bell-examples}
 e_2&=\tfrac12(p_1^2-p_2),\notag\\
 e_3&=\tfrac16(p_1^3-3p_1p_2+2p_3),\notag\\
 e_4&=\tfrac1{24}
       (p_1^4-6p_1^2p_2+3p_2^2+8p_1p_3-6p_4).
\end{align}

\subsection{Explicit Gamma and Stieltjes identities}

For zero-centered first derivatives define
\[
 L(y,z)=\log\Gamma(z+1)-\log\Gamma(y),\qquad
 R_j(y,z)=\zeta^{[j]}(0,z+1)-\zeta^{[j]}(0,y).
\]
The classical identity
\(\zeta^{[1]}(0,x)=\log\Gamma(x)-\tfrac12\log(2\pi)\)
\cite{DLMF} makes \(p_1=L\), while \(p_j=R_j\) for
\(j\ge2\).  The source's first two-marker target is therefore
\begin{equation}\label{gm:gamma-square}
 \boxed{\left.\partial_u\partial_v
       \mathcal H_{a,b}(s,u,v,t)\right|_{u=v=0}
       =\frac12\mathcal E_{a,b}^{s,t}[L^2-R_2].}
\end{equation}
This is an identity at arbitrary independent \(s,t\) and arbitrary
right-half-plane endpoints.  It supplies ordinary absolutely convergent
sums through \eqref{gm:ordinary-E} whenever \(\Re s,\Re t>3\):
\(\log\Gamma(x)=O(x\log x)\) and
\(\zeta^{[2]}(0,x)=O(x\log^2x)\) on each fixed lattice give
summability of every separated letter.  Elsewhere it is the specified
entire continuation.  The third and fourth identical marked clusters
are obtained from \eqref{gm:Bell-examples}; no new ordered coordinate
is required for these symmetric identical letters.

To connect this directly with normalized antiderivatives, write the
Hurwitz parameter-derivative law in its entire form
\[
 \partial_x\zeta(u,x)=-1-uQ(u+1,x).
\]
Differentiating at \(u=0\), for every \(j\ge1\), gives
\begin{equation}\label{gm:primitive}
 \partial_x\zeta^{[j]}(0,x)=(-1)^j j\gamma_{j-1}(x),\qquad
 R_j(y,z)=(-1)^j j\int_y^{z+1}\gamma_{j-1}(x)\,dx.
\end{equation}
The paths lie in the right half-plane; the constants are fixed by
vanishing when \(z+1=y\).  Thus \eqref{gm:gamma-square} is also
\begin{equation}\label{gm:gamma-integral}
 \left.\partial_u\partial_v\mathcal H_{a,b}(s,u,v,t)\right|_{0,0}
  =\frac12\mathcal E_{a,b}^{s,t}\!\left[
       L(y,z)^2-2\int_y^{z+1}\gamma_1(x)\,dx\right].
\end{equation}
The logarithm in \(L\) remains the difference of analytic log-Gamma
branches; it is not replaced by an independently chosen logarithm of
a ratio.

An example in which the merging itself creates a Stieltjes coefficient
uses two half-order marks.  Set
\(A(y,z)=\zeta^{[1]}(1/2,z+1)-\zeta^{[1]}(1/2,y)\).
Then
\begin{equation}\label{gm:half-pair}
 \left.\partial_u\partial_v\mathcal H_{a,b}(s,u,v,t)
                                      \right|_{u=v=1/2}
 =\frac12\mathcal E_{a,b}^{s,t}
          [A(y,z)^2-\gamma_2(z+1)+\gamma_2(y)].
\end{equation}
Indeed the merged index is one and its total derivative order is two.

For an unequal pair, the symmetric distinction is essential.  Define
\[
 J_{0,1}(s,t)=\left.\partial_u\partial_v
 \bigl\{\mathcal H_{a,b}(s,u,v,t)+\mathcal H_{a,b}(s,v,u,t)\bigr\}
                                                    \right|_{u=0,v=1}.
\]
The exact mixed Gamma--Stieltjes formula is
\begin{equation}\label{gm:mixed-pair}
 J_{0,1}(s,t)=\mathcal E_{a,b}^{s,t}\!\left[
 -L(y,z)\{\gamma_1(z+1)-\gamma_1(y)\}
                         -\gamma_2(z+1)+\gamma_2(y)\right].
\end{equation}
This determines the sum of the two ordered derivatives, not each
summand separately.

\subsection{Several gaps and the exact scope of the reduction}

For surviving variables \(x_1,\ldots,x_d\), the three types of
completed gaps are
\begin{equation}\label{gm:all-gaps}
 \begin{cases}
 Q(u,x_1+1)-Q(u,a),&\text{above the first variable},\\
 Q(u,x_{j+1}+1)-Q(u,x_j),&\text{between successive variables},\\
 Q(u,b)-Q(u,x_d),&\text{below the last variable}.
 \end{cases}
\end{equation}
On an absolute-convergence domain, $Z_c$ is a quasi-shuffle character
on lattice functions with pointwise multiplication. Characters with
values in a commutative algebra form a group under deconcatenation
convolution; therefore $Z_a^{-1}*Z_b$ has this property for decorated
letters as well. The continuation proved above preserves these
identities. The upper boundary in \eqref{gm:all-gaps} follows from the
one-tail lemma at the first position and the shift law. The lower
boundary follows by subtracting that lemma at the last position from
the singleton stuffle identity, using
$\mathcal H_{a,b}(u)=Q(u,b)-Q(u,a)$.

Apply the partition polynomial independently in each gap and multiply
the resulting polynomials before expanding into decorated letters.
The proof by insertion applies gap by gap, and
Theorem~\ref{gm:closure} gives joint continuation, including all shared
spectral resonances.  It determines the ordered word whenever each gap
contains one mark or an identical marked cluster.  With no surviving
variables, the sole gap is \(Q(u,b)-Q(u,a)\), and the construction
reduces to the classical symmetric-sum formula for the relative
interpolant.

The accompanying program verifies the local symmetric identity through
five independent labels by exact integer arithmetic on formal decorated
words.  All endpoint factors and shorter words cancel, leaving precisely
the \(k!\) permutations.  A separate numerical check uses the original
nested \(Z_c\)-series and relative quotient at nonintegral complex
endpoints: increasing the cutoff from 80 to 160 terms reduces its
absolute residual from \(1.66\times10^{-36}\) to
\(2.69\times10^{-40}\).  An independent reconstruction of the actual
zero-marker derivative at those noninteger endpoints, using convergent
nested sums and a symmetrized top-tail continuation, gives residuals
\(1.40\times10^{-27}\) and \(3.73\times10^{-31}\) at the same
cutoffs.  Individual singular double-zeta terms are never differentiated
at \((0,0)\).  Five direct finite-interval logarithmic
comparisons test negative centers, the Gamma square, three marks, and
merged Stieltjes resonances at 60 decimal digits; each residual is below
\(8\times10^{-62}\).  These are diagnostics of the algebra and
normalization, not interval certificates or substitutes for the
analytic proof.

The audited source correctly withheld the several-marker claim; no
correction to its single-marker theorem is needed.  The present theorem
answers its first Q4 target and the fully symmetric part of Q5.  A
remaining precise question is to represent the antisymmetric ordered
two-marker component with unequal marks.  Another is to classify which
unmarked nonpositive deletions preserve this symmetric span: such
deletions can shift different marked centers unequally, so their full
ordered extension is not implied by the Bell formula.

% END INCLUDED FILE: sections/04_harmonic_gaps.tex


% BEGIN INCLUDED FILE: sections/05_spectral_faces.tex
% Self-contained section fragment; requires amsmath, amssymb, amsthm.
% The parent document supplies theorem, lemma, proposition, corollary, remark.
% Labels and bibliography keys in this contribution have the prefix sd:.
\section{Higher spectral derivatives through the faces of a simplex}
\label{sd:sec:main}

The incoming beta-polarization report asks for explicit higher-order
subtractions on a Dirichlet simplex and for their extension to odd
polynomial weights \cite[Further research, items 3--4]{sd:HigherDefects}.
The incoming quartic report already gives a finite-subtraction compiler
at every spectral order \cite[Section 5]{RepoQuartic}.  We give a different
answer: a finite representation by ordinary simplex integrals of
\emph{lower} one-mass derivatives.  Its useful structural feature is that
an interaction involving a face with $k$ vertices cannot appear before
the $k$th spectral derivative, after an explicitly computable zero-jet
polynomial has been separated.

The scalar Dirichlet parameter identity is classical.  It is already
used for mixed twists in the canonical manuscript
\cite[\texttt{07-mixed-twist-identities.tex}]{RepoMain}; logarithmic
Dirichlet averages have a substantial classical theory
\cite{sd:Carlson1987,sd:Carlson1991}.  The contribution here is the explicit
face subtraction, its absolute-convergence proof, the resulting
all-order formula, and the rational zero-jet term.  This is a claim of
new content relative to the inspected repository, not an assertion of
exhaustive literature priority.

\subsection{Spectral normalization and the zero-jet polynomial}
\label{sd:subsec:setup}

Fix $a>0$, $r\ge2$, real masses $B_1,\ldots,B_r\ge0$, and a polynomial
$W(x)=\sum_{d=0}^{D}w_dx^d$.  Write $x_n=n+a$ and initially define
\begin{align}
 E_W(s;B)&=\sum_{n\ge0}W(x_n)(x_n^2+B)^{-s},
       \label{sd:eq:one}\\
 Z_{W,\boldsymbol B}(s)&=\sum_{n\ge0}W(x_n)
                    \prod_{i=1}^r(x_n^2+B_i)^{-s},\label{sd:eq:many}\\
 \mathcal D_W(s)&=\frac1r\sum_{i=1}^rE_W(rs;B_i)
                          -Z_{W,\boldsymbol B}(s),\qquad
 \Delta_m(W)=\mathcal D_W^{(m)}(0).
       \label{sd:eq:defect}
\end{align}
The sums start in their half-planes of absolute convergence and are then
continued meromorphically.  All logarithms of the positive quadratic
factors are real.  Let
\begin{equation}\label{sd:eq:jets}
 F_j(B)=\left.\partial_s^j E_W(s;B)\right|_{s=0}.
\end{equation}

\begin{lemma}[The polynomial at spectral order zero]
\label{sd:lem:zero}
The spectral functions in \eqref{sd:eq:one}--\eqref{sd:eq:many} are
holomorphic at zero.  The one-mass zero jet is
\begin{equation}\label{sd:eq:zero}
 F_0(B)=\sum_{d=0}^{D}w_d\zeta(-d,a)
           +\sum_{h\ge1}p_h B^h,\qquad
 p_h=\frac{(-1)^h}{2h}\,w_{2h-1},
\end{equation}
where missing polynomial coefficients are zero.  In particular $F_0$
is mass-independent for even $W$.
\end{lemma}
\begin{proof}
For sufficiently large $x$, expand the quadratic factor with a finite
remainder:
\[
 (x^2+B)^{-s}=x^{-2s}\sum_{h=0}^{K}
                 \frac{(-1)^h(s)_h}{h!}B^h x^{-2h}
                  +\text{remainder}.
\]
After multiplication by $x^d$, the removed terms sum to
$(-1)^h(s)_hB^h\zeta(2s+2h-d,a)/h!$.  Choose $K$ large enough
that the remainder and its required derivatives converge normally near
zero.  A finite head is retained exactly, so no restriction such as
$B<a^2$ is needed.  For $h\ge1$, the coefficient vanishes at zero
unless $2h-d=1$.  In that case
\[
 \frac{(s)_h}{h!}=\frac{s}{h}+O(s^2),\qquad
 \zeta(1+2s,a)=\frac1{2s}+O(1),
\]
which gives the term $(-1)^h B^h/(2h)$ in \eqref{sd:eq:zero}.
The same subtraction for the product of quadratic factors has positive
Taylor-degree coefficients vanishing at $s=0$; they cancel its possible
simple Hurwitz poles.  This also proves holomorphy of the product
spectrum.  The argument is locally uniform in the masses and in $a>0$.
\end{proof}

\subsection{An explicit subtraction on every face}
\label{sd:subsec:faces}

For a nonempty set $S\subseteq\{1,\ldots,r\}$, let
\[
 \Sigma_S=\{(t_i)_{i\in S}:t_i\ge0,\ \sum_{i\in S}t_i=1\}.
\]
The measure is ordinary $(|S|-1)$-dimensional coordinate measure:
$dt_1\cdots dt_{k-1}$ after the final coordinate has been eliminated.
Thus the simplex has volume $1/(k-1)!$, where $k=|S|$.
For a nonempty $T\subseteq S$ put
\[
 t_T=\sum_{i\in T}t_i,\qquad
 B_T=\frac{\sum_{i\in T}t_i B_i}{t_T},\qquad
 G_T(F;t)=t_TF(B_T).
\]
Set $G_T=0$ when $t_T=0$.  Define the face subtraction
\begin{equation}\label{sd:eq:H}
 \boxed{\displaystyle
 \mathcal H_S(F;t)=
  \sum_{\varnothing\ne T\subseteq S}
       (-1)^{|S|-|T|}\,t_TF(B_T).}
\end{equation}
For a pair this is the difference from the affine chord:
\begin{equation}\label{sd:eq:Hpair}
 \mathcal H_{ij}(F;u)
   =F(uB_i+(1-u)B_j)-uF(B_i)-(1-u)F(B_j).
\end{equation}
For three vertices it is
\begin{align}\label{sd:eq:Htriple}
 \mathcal H_{ijk}(F;t)={}&F(t_iB_i+t_jB_j+t_kB_k)
                  +\sum_{\ell\in\{i,j,k\}}t_\ell F(B_\ell)\notag\\
 &-\sum_{\{\ell,m\}\subseteq\{i,j,k\}}
   (t_\ell+t_m)F\!\left(
            \frac{t_\ell B_\ell+t_mB_m}{t_\ell+t_m}\right).
\end{align}
These formulas remain regular when masses coincide.  They use neither
division by a mass difference nor an unspecified finite part.

\begin{lemma}[Boundary cancellation and ordinary integrability]
\label{sd:lem:boundary}
If $F\in C^1([\min B_i,\max B_i])$, then $\mathcal H_S$ vanishes
on every proper coordinate face and
\begin{equation}\label{sd:eq:bound}
 |\mathcal H_S(F;t)|\le C\min_{i\in S}t_i.
\end{equation}
For $k=|S|\ge2$, the integral
\begin{equation}\label{sd:eq:integralhol}
 \int_{\Sigma_S}\mathcal H_S(F;t)
                            \prod_{i\in S}t_i^{s-1}\,dt
\end{equation}
is absolutely convergent for $\operatorname{Re}s>-1/(k-1)$.
It is holomorphic there, with locally normally convergent fixed
logarithmic derivatives.  The same conclusions hold for a holomorphic
family of $F$ whose $C^1$ norms are locally bounded.
\end{lemma}
\begin{proof}
The perspective $G_T=t_TF(B_T)$ is Lipschitz on the nonnegative
coordinate cone.  Away from $t_T=0$,
\[
 \partial_{t_i}G_T=F(B_T)+F'(B_T)(B_i-B_T)\quad(i\in T),
\]
so its coordinate derivatives have a uniform bound.  Continuity at the
origin extends the Lipschitz bound there.  If $t_i=0$, the terms indexed
by $T$ and $T\cup\{i\}$ cancel in \eqref{sd:eq:H}.  Changing that
coordinate from zero to $t_i$ therefore gives a bound $Ct_i$ for
every $i$, proving \eqref{sd:eq:bound}.

Cover the simplex by the $k$ regions where a specified coordinate is
at least $1/k$.  On such a region,
\[
 \min_i t_i\le
        \left(\prod_{i\ne i_{\max}}t_i\right)^{1/(k-1)}.
\]
The remaining coordinate factors are dominated by a product of
integrable powers whenever $\operatorname{Re}s+1/(k-1)>0$.
Any fixed power of $\sum_i|\log t_i|$ remains integrable on smaller
closed half-planes.  Dominated holomorphic parameter integration proves
all assertions.  In particular, every integral below with the weight
$\prod_i t_i^{-1}$ is an ordinary absolutely convergent integral.
\end{proof}

\begin{theorem}[Exact decomposition by simplex faces]
\label{sd:thm:master}
Put
\begin{equation}\label{sd:eq:Q}
 Q_k(s)=\frac{\Gamma(1+ks)}{\Gamma(1+s)^k}.
\end{equation}
The following identity holds meromorphically for
$\operatorname{Re}s>-1/(r-1)$, with ordinary face integrals away
from spectral poles, and in particular as a holomorphic germ at zero:
\begin{equation}\label{sd:eq:master}
 \boxed{\displaystyle
 \mathcal D_W(s)=-\frac1r\sum_{k=2}^{r}s^{k-1}Q_k(s)
   \sum_{|S|=k}\int_{\Sigma_S}
     \mathcal H_S(E_W(rs;\cdot);t)
                         \prod_{i\in S}t_i^{s-1}\,dt.}
\end{equation}
\end{theorem}
\begin{proof}
The product of Gamma integrals, with variables changed to a radial
coordinate and simplex proportions, gives
\[
 \prod_{i=1}^r y_i^{-s}
 =\frac{\Gamma(rs)}{\Gamma(s)^r}
   \int_{\Sigma_{\{1,\ldots,r\}}}
        \left(\sum_i t_i y_i\right)^{-rs}
                            \prod_i t_i^{s-1}\,dt
\]
for $y_i>0$ and $\operatorname{Re}s>0$.  Apply this to
$y_i=x_n^2+B_i$ and sum where the spectra converge absolutely.
Finite subset inversion gives
\[
 F\!\left(\sum_i t_iB_i\right)
       =\sum_{\varnothing\ne S\subseteq\{1,\ldots,r\}}
                    \mathcal H_S(F;t).
\]
On the full simplex, each $\mathcal H_S$ is homogeneous of degree one
in the coordinates belonging to $S$.  Under the normalized symmetric
Dirichlet weight, $q=t_S$ has mean $|S|/r$ and is independent of the
normalized proportions $t_i/q$, $i\in S$.  Equivalently, this follows
by integrating first in $q$ and the complementary proportions using
the ordinary beta integral.  Therefore, for $k=|S|$,
\[
 \mathbb E_{r,s}\mathcal H_S(F;t)
 =\frac{k}{r}\frac{\Gamma(ks)}{\Gamma(s)^k}
       \int_{\Sigma_S}\mathcal H_S(F;t)
                                 \prod_{i\in S}t_i^{s-1}\,dt.
\]
The singleton terms are $r^{-1}\sum_iF(B_i)$.  Subtracting them and
using
$\frac{k}{r}\Gamma(ks)/\Gamma(s)^k=s^{k-1}Q_k(s)/r$
proves \eqref{sd:eq:master} in the initial half-plane.
The spectral finite subtractions in Lemma~\ref{sd:lem:zero}, together
with Lemma~\ref{sd:lem:boundary}, give meromorphic continuation of the
integrals and the asserted holomorphic germs.  The equality follows
there by continuation from the common convergence domain.
\end{proof}

\subsection{The entire zero-jet contribution is a finite rational function}
\label{sd:subsec:local}

Separate $E_W(rs;B)=F_0(B)+(E_W(rs;B)-F_0(B))$ in
\eqref{sd:eq:master}.  Only the mass polynomial in
\eqref{sd:eq:zero} contributes to the first part.

\begin{theorem}[Rational zero-jet contribution]
\label{sd:thm:local}
The part of $\mathcal D_W(s)$ arising from $F_0$ is
\begin{equation}\label{sd:eq:local}
 \boxed{\displaystyle
 \mathcal L_W(s)=\sum_{h\ge1}p_h\left\{
 \frac1r\sum_{i=1}^r B_i^h-
 \frac{h!}{(rs)_h}
       \sum_{h_1+\cdots+h_r=h}
            \prod_{i=1}^r\frac{(s)_{h_i}B_i^{h_i}}{h_i!}
                                      \right\}.}
\end{equation}
All apparent singularities at $s=0$ are removable, and
$\mathcal L_W(0)=0$.  The term $h=1$ vanishes identically.  For
$m\ge1$, set
\begin{equation}\label{sd:eq:Ccoeff}
 C_{h,j}=[z^h]\prod_{i=1}^r(1-B_i z)^{j/r},
             \qquad 1\le j\le h-1.
\end{equation}
Then
\begin{equation}\label{sd:eq:localderivative}
 \boxed{\displaystyle
 \mathcal L_W^{(m)}(0)=
 \sum_{h\ge2}p_h(-1)^{m+1}m!r^m h!
  \sum_{j=1}^{h-1}
    \frac{(-1)^jC_{h,j}}{j!(h-1-j)!\,j^{m+1}}.}
\end{equation}
The coefficient extraction in \eqref{sd:eq:Ccoeff} is finite and
algebraic; it makes no branch choice beyond the formal series with
constant coefficient one.
\end{theorem}
\begin{proof}
Dirichlet moments give
\[
 \mathbb E_{r,s}\left(\sum_i t_iB_i\right)^h
  =\frac{h!}{(rs)_h}
       \sum_{h_1+\cdots+h_r=h}
                    \prod_i\frac{(s)_{h_i}B_i^{h_i}}{h_i!}.
\]
This proves \eqref{sd:eq:local}.  At $s=0$ its continuation is the
average of the vertex values: only compositions supported on a single
index survive after the common factor $s$ has been cancelled.

For the derivative formula put $v=rs$.  The $h$th moment becomes
\[
 M_h(v)=\frac{h!}{(v)_h}
                   [z^h]\prod_i(1-B_i z)^{-v/r}.
\]
After cancellation at zero, it is a rational function of numerator and
denominator degree at most $h-1$.  Its polynomial part is constant.
Its possible poles are simple and lie at $v=-j$, $1\le j<h$, with
residues
\[
 \operatorname{Res}_{v=-j}M_h(v)
   =\frac{(-1)^j h!C_{h,j}}{j!(h-1-j)!}.
\]
Differentiate the partial fractions $m$ times at $v=rs=0$ and subtract
from the vertex average.  This is exactly
\eqref{sd:eq:localderivative}.
\end{proof}

For example, let $V(\boldsymbol B)=\sum_{i<j}(B_i-B_j)^2$.  For the
first nontrivial odd weight,
\begin{equation}\label{sd:eq:cubiclocal}
 \mathcal L_{x^3}(s)=\frac{sV(\boldsymbol B)}{4r(1+rs)},\qquad
 \mathcal L_{x^3}^{(m)}(0)
       =\frac{m!(-r)^{m-1}}{4r}V(\boldsymbol B).
\end{equation}
More generally, after division by $m!$, the sequence in
\eqref{sd:eq:localderivative} is a finite linear combination of powers
$(-r/j)^m$.  It therefore satisfies an explicit finite recurrence in
$m$.  This rational contribution is the part supplied by the zero
one-mass jet.  It should not be read as saying that Stieltjes or residue
terms inside the lower jets $F_1,F_2,\ldots$ disappear.

\subsection{Every derivative uses only finitely many lower jets}
\label{sd:subsec:allorders}

For $k\ge2$ and $q\ge0$ define
\begin{equation}\label{sd:eq:P}
 P_{k,q}(L)=[z^q]\exp\left(
 Lz+\sum_{n\ge2}\frac{(-1)^n\zeta(n)}n(k^n-k)z^n\right).
\end{equation}
Only finitely many terms of the exponential are required at any fixed
$q$.  In particular,
\[
 P_{k,0}=1,\qquad P_{k,1}=L,\qquad
 P_{k,2}=\frac{L^2}{2}+\frac{k^2-k}{2}\zeta(2).
\]
Put $\ell_S(t)=\sum_{i\in S}\log t_i$.

\begin{theorem}[All-order formula and face-size restriction]
\label{sd:thm:allorders}
For every integer $m\ge1$,
\begin{equation}\label{sd:eq:allorders}
 \boxed{\begin{aligned}
 \Delta_m(W)={}&\mathcal L_W^{(m)}(0)\\
 &-\frac{m!}{r}
  \sum_{k=2}^{\min(r,m)}\ \sum_{|S|=k}
  \sum_{j=1}^{m-k+1}\frac{r^j}{j!}
   \int_{\Sigma_S}\mathcal H_S(F_j;t)
       P_{k,m-k+1-j}(\ell_S(t))
                     \frac{dt}{\prod_{i\in S}t_i}.
 \end{aligned}}
\end{equation}
Every displayed simplex integral is absolutely convergent.  After the
explicit zero-jet term is removed, a $k$-vertex face starts at order
$s^k$ and uses one-mass derivatives of order at most $m-k+1$ in the
$m$th derivative.
\end{theorem}
\begin{proof}
The classical Taylor series of $\log\Gamma(1+z)$ gives
\[
 Q_k(s)=\exp\left(
       \sum_{n\ge2}\frac{(-1)^n\zeta(n)}n(k^n-k)s^n\right);
\]
the Euler-constant terms cancel.  Once $F_0$ has been removed, the
inner function in \eqref{sd:eq:master} is
$\sum_{j\ge1}r^js^jF_j/j!$.  Multiplication by $s^{k-1}$ proves
that this face starts at degree $k$.  Extract the coefficient of $s^m$
from \eqref{sd:eq:master}; the remaining exponential is precisely
\eqref{sd:eq:P}.  Lemma~\ref{sd:lem:boundary} justifies extraction under
each integral.  Multiplication by $m!$ gives the theorem.
\end{proof}

For a compact statement of the first instances write
\[
 I_S(F)=\int_{\Sigma_S}\mathcal H_S(F;t)
                       \frac{dt}{\prod_i t_i},\qquad
 J_S(F)=\int_{\Sigma_S}\mathcal H_S(F;t)\ell_S(t)
                       \frac{dt}{\prod_i t_i},
 \qquad L_m=\mathcal L_W^{(m)}(0).
\]
All subset sums below use distinct labels.  The complete second and
third derivatives are
\begin{align}
 \Delta_2(W)&=L_2-2\sum_{|S|=2}I_S(F_1),
                    \label{sd:eq:second}\\
 \Delta_3(W)&=L_3-3r\sum_{|S|=2}I_S(F_2)
                  -6\sum_{|S|=2}J_S(F_1)
                  -6\sum_{|S|=3}I_S(F_1).
                    \label{sd:eq:third}
\end{align}
The fourth derivative is
\begin{align}\label{sd:eq:fourth}
 \Delta_4(W)={}&L_4-4r^2\sum_{|S|=2}I_S(F_3)
                    -12r\sum_{|S|=2}J_S(F_2)\notag\\
 &-24\sum_{|S|=2}\int_{\Sigma_S}\mathcal H_S(F_1;t)
             \left(\frac{\ell_S(t)^2}{2}+\zeta(2)\right)
                                  \frac{dt}{\prod_i t_i}\notag\\
 &-12r\sum_{|S|=3}I_S(F_2)-24\sum_{|S|=3}J_S(F_1)
                              -24\sum_{|S|=4}I_S(F_1).
\end{align}
A subset sum is empty when its size exceeds $r$.  The term $\zeta(2)$
in the pair integral of \eqref{sd:eq:fourth} comes from $Q_2$ and is
part of the normalization.  In particular the third derivative needs
only pair integrals of the second one-mass jet and pair and triple
integrals of the first jet.

\subsection{Explicit Gamma integrands and a nonzero face family}
\label{sd:subsec:gamma}

For $W=1$, first-derivative factorization gives the classical expression
\begin{equation}\label{sd:eq:gamma}
 F_1(B)=\log\Gamma(a+i\sqrt B)+\log\Gamma(a-i\sqrt B)
                                      -\log(2\pi),
\end{equation}
using the analytic logarithm of Gamma in the right half-plane.
The corresponding even-weight expression, inherited from
first-derivative factorization, is
\begin{equation}\label{sd:eq:evenhurwitz}
 F_1(B)=\sum_{\varepsilon=\pm1}\sum_{j=0}^{2d}
   \binom{2d}{j}(-\varepsilon i\sqrt B)^{2d-j}
               \zeta'(-j,a+\varepsilon i\sqrt B),
             \qquad W=x^{2d}.
\end{equation}
Thus the largest face size $k=m$ in \eqref{sd:eq:allorders} requires
only ordinary Gamma at $W=1$, and a finite Hurwitz-derivative expression
for every even polynomial weight.  Its coefficients and integration
kernels are fully explicit.

The following identity shows that the higher faces are substantive
terms of this particular representation.

\begin{proposition}[A factorial identity for every face size]
\label{sd:prop:quadraticface}
For $k=|S|\ge2$,
\begin{equation}\label{sd:eq:quadraticface}
 \boxed{\displaystyle
 I_S(B\mapsto B^2)=(-1)^{k-1}(k-2)!
                             \sum_{\{i,j\}\subseteq S}(B_i-B_j)^2.}
\end{equation}
Consequently, if $B_i=B+\epsilon\beta_i$ and $F$ is $C^3$ near $B$,
\begin{equation}\label{sd:eq:collisionleading}
 I_S(F)=\frac{\epsilon^2F''(B)}2(-1)^{k-1}(k-2)!
           \sum_{\{i,j\}\subseteq S}(\beta_i-\beta_j)^2
                         +O(\epsilon^3).
\end{equation}
\end{proposition}
\begin{proof}
On the nonnegative coordinate cone the quadratic perspective has the
variance decomposition
\[
 t_T B_T^2=\sum_{i\in T}t_iB_i^2
       -\sum_{\{i,j\}\subseteq T}
                 \frac{t_it_j}{t_T}(B_i-B_j)^2.
\]
The first sum is annihilated by the face subtraction.  Writing
$t_T^{-1}=\int_0^\infty e^{-vt_T}\,dv$ in the second gives
\begin{align*}
 \mathcal H_S(B^2;t)
  ={}&(-1)^{k-1}\sum_{\{i,j\}\subseteq S}(B_i-B_j)^2t_it_j\\
 &\qquad\times\int_0^\infty e^{-v(t_i+t_j)}
                 \prod_{\ell\in S\setminus\{i,j\}}
                         (1-e^{-vt_\ell})\,dv.
\end{align*}
Since $\mathcal H_S$ is homogeneous of degree one, radial integration
shows that $I_S$ equals the integral of
$e^{-\sum_i t_i}\mathcal H_S(t)/\prod_i t_i$ on the positive
orthant.  The radial factor is $\int_0^\infty e^{-\rho}\,d\rho=1$.
After the displayed sign is removed, every term is nonnegative, so
Tonelli's theorem applies.  The two distinguished coordinates give
$(1+v)^{-2}$, and each other coordinate gives
\[
 \int_0^\infty e^{-t}\frac{1-e^{-vt}}{t}\,dt=\log(1+v).
\]
The remaining integral is
\[
 \int_0^\infty\frac{\log^{k-2}(1+v)}{(1+v)^2}\,dv=(k-2)!.
\]
This proves \eqref{sd:eq:quadraticface}.  Taylor-expand $F(B+\epsilon
\beta)$ through degree two.  The constant and affine parts are
annihilated; Lemma~\ref{sd:lem:boundary} controls the integrated
remainder uniformly, proving \eqref{sd:eq:collisionleading}.
\end{proof}

For \eqref{sd:eq:gamma},
\[
 F_1''(B)=\sum_{n\ge0}\frac1{(x_n^2+B)^2}>0.
\]
Hence every face size has a nonzero quadratic collision coefficient
when the perturbations are not all equal.  This does not prove
irreducibility into any other chosen function class.  In particular,
$\mathcal H_S$ need not vanish merely because two of several masses
coincide; ``face contribution'' is the precise description used here.

\subsection{Periodic weights and cancellation of the mean}
\label{sd:subsec:periodic}

\begin{corollary}[Mean-zero periodic weights remove the zero-jet term]
\label{sd:cor:periodic}
Let $c_n$ be $q$-periodic, and insert $c_n$ in every spectral sum.
Theorems~\ref{sd:thm:master} and \ref{sd:thm:allorders} still hold.
The rational zero-jet contribution is
\begin{equation}\label{sd:eq:periodiclocal}
 \mathcal L_{c,W}(s)=\overline c\,\mathcal L_W(s),\qquad
                  \overline c=\frac1q\sum_{v=0}^{q-1}c_v.
\end{equation}
Thus $\overline c=0$ removes this term for \emph{every} polynomial
weight, including odd weights.
\end{corollary}
\begin{proof}
The auxiliary Dirichlet series is
\[
 \sum_{n\ge0}c_n(n+a)^{-z}
   =q^{-z}\sum_{v=0}^{q-1}c_v
                           \zeta\!\left(z,\frac{a+v}{q}\right).
\]
Its residue at $z=1$ is $\overline c$.  Repeat the proof of
Lemma~\ref{sd:lem:zero}; the zero jet becomes
\[
 F_{c,0}(B)=\sum_d w_dq^d
       \sum_v c_v\zeta\!\left(-d,\frac{a+v}{q}\right)
                    +\overline c\sum_{h\ge1}p_hB^h.
\]
The mass-independent part is annihilated by every face subtraction.
All other steps use linearity and normal convergence; periodic weights
are bounded and cause no additional convergence issue.
\end{proof}

For $W=1$ the first jet has the explicit finite dictionary
\begin{align}\label{sd:eq:periodicgamma}
 F_{c,1}(B)=\sum_{v=0}^{q-1}c_v\bigg\{&
 \log\Gamma\!\left(\frac{a+v+i\sqrt B}{q}\right)
 +\log\Gamma\!\left(\frac{a+v-i\sqrt B}{q}\right)\notag\\
 &-\log(2\pi)-2\log q
                         \left(\frac12-\frac{a+v}{q}\right)\bigg\}.
\end{align}
This follows by scaling each residue class to
$q^{-2s}E_1(s;B/q^2)$ with shift $(a+v)/q$, where the subscript
$1$ here denotes the polynomial weight.  Formula
\eqref{sd:eq:periodicgamma} therefore supplies ordinary Gamma
integrands for alternating weights and character projections.  The
mass-independent constants can be omitted inside $\mathcal H_S$,
but are retained in the full first-jet identity.

\subsection{Independent evaluation and research consequences}
\label{sd:subsec:validation}

For $a=2$, $W=1$, $r=3$, and masses
$(B_1,B_2,B_3)=(0.2,0.7,1.3)$, the directly evaluated spectral derivative is
\begin{equation}\label{sd:eq:numeric}
 \Delta_3=-1.8999673234892670916122436912924709222747240234\ldots.
\end{equation}
One calculation used a finite ordinary spectral sum and an independently
summed binomial--Hurwitz tail.  Changing the head length and tail degree
from $(55,13)$ to $(85,17)$ changed the result by $3.04\times10^{-46}$.
A second calculation evaluated \eqref{sd:eq:third}: the pair terms used
convergent one-mass expansions, while the triple integral used Gamma
values from \eqref{sd:eq:gamma}.  The errors in the full face formula
at Gauss--Legendre orders $16$, $24$, and $32$, with $45$-digit
arithmetic, were respectively
$1.62\times10^{-18}$, $2.54\times10^{-27}$, and
$8.20\times10^{-32}$.  The last error is limited by the retained
one-mass expansion.  The triple integral was
\[
 I_{\{1,2,3\}}(F_1)
    =0.0564600947235384049461537670929012515\ldots.
\]
For stable numerical integration the simplex was divided into its six
coordinate orderings.  On each, use
$t=(1,u,uv)/(1+u+uv)$; the weighted measure becomes
$dt/(t_1t_2t_3)=du\,dv/(uv)$, and the face cancellation removes the
corner singularity.  The code and precision-dependent outputs are
included with the article.  These are numerical diagnostics, not
interval enclosures or replacements for the proofs.

The results close the explicit face-subtraction problem and supply a
complete odd-weight zero-jet correction in this lower-jet framework.
Further concrete questions are to evaluate special Gamma face integrals
in a fixed polylogarithmic function class; to obtain comparable
multivariable formulas when the spectral orders of the factors vary
independently; and to derive closed all-order collision coefficients for
the individual face integrals beyond \eqref{sd:eq:collisionleading}.
The theorem itself does not require such reductions and makes no claim
that the unresolved Gaussian period conjectures are settled.

% END INCLUDED FILE: sections/05_spectral_faces.tex


% BEGIN INCLUDED FILE: sections/06_audit_and_research.tex
\section{Source audit, exact verification, and integration decisions}
\label{sec:audit}

\subsection{What changes in the source status}

The article supplies advances at specific research interfaces. It does
not replace the source's definitions or silently promote other open
claims. The following status changes are proposed for integration.

\begingroup
\small
\begin{longtable}{@{}p{0.27\textwidth}p{0.65\textwidth}@{}}
\toprule
Source target & Proposed status after review of this article\\
\midrule
\endhead
Quartic report, R4: fifth- and higher-power cancellation &
The universal recurrence and anchored generator are proved in
Theorems~\ref{tp:thm:recurrence} and \ref{tp:thm:generator}. Exact
eight- and eleven-coordinate rows are supplied for powers five and six.
Theorem~\ref{et:main} treats every repeated even order. The minimal
general coordinate count remains a separate question.\\[5pt]
Quartic report, R5: resonant parity &
The complete Laurent normalization and all coefficients are proved in
Theorem~\ref{tp:thm:Laurent}. The local spectral Taylor series has exact
radius two. A uniform finite arithmetic reduction as the resonance
index varies is not proved.\\[5pt]
Quartic report, R1: smaller order ten &
The complete all-$b$ statement for $b\ge6$ is proved in
Theorem~\ref{mt10:main}; the new missing certificate is $b=962$.
The smaller-order-twelve and unrestricted unequal-order questions remain
open here.\\[5pt]
Harmonic-gap continuation, Q4 and Q5 &
The first two-marker target and all identical marked clusters are
proved by \eqref{gm:gamma-square} and \eqref{gm:identical}.
Corollary~\ref{gm:marked} treats the fully symmetric unequal-marker
component. General individual ordered unequal-marker reductions remain
open.\\[5pt]
Higher-defect report, face and odd-weight questions &
Theorems~\ref{sd:thm:master}--\ref{sd:thm:allorders} give the ordinary
face decomposition and its full rational zero-jet term. The periodic
extension is Corollary~\ref{sd:cor:periodic}. Special-function evaluation
of general higher face integrals remains open.\\
\bottomrule
\end{longtable}
\endgroup

The general all-power generator is not the already known finite-span
theorem for powers of the ordinary harmonic number $H_n$. Here the
objects are powers of the decaying polygamma tail $\zeta(2,n+1)$.
Likewise the face theorem is not a new claim that a Dirichlet mixture
exists: it is the explicit decomposition and integrability of its
boundary-subtracted pieces. The relative Hurwitz interpolation itself
is inherited from published work. These separations prevent duplicate
claims during integration.

\subsection{Corrections and normalization details}

No new false theorem was found in the targeted rank and harmonic-gap
material reviewed for this continuation. In particular, the source's
single-marker theorem correctly left the several-marker extension open.
The lower-index-zero binomial omission in the reconstruction matrices
was already explicit in the canonical fifth-order report and in the
incoming extensions. It must be retained as an existing convention,
not reported as a newly discovered erratum.

The following details are essential when the new results are inserted.
They are precise safeguards for the supplied proof chain.

\begin{enumerate}
\item In the centered finite-cutoff argument use $n=0,\ldots,N-1$.
The even-power subtraction sums then have zero constant coefficient.
Changing to a cutoff ending at $N$ requires its endpoint polynomial.
The missing $n=0$ endpoint also produces the exceptional
$-d\delta_{k,2}\zeta(2)^{d-1}/2$ in
\eqref{tp:eq:recurrence}.

\item The tail expansions are asymptotic. Only a finite subtraction,
with a normally convergent remainder, is used at any specified
spectral point. A fixed infinite asymptotic series is not summed as
an ordinary power series in $1/n$.

\item The words ``eight coordinates,'' ``eleven coordinates,'' and
``rank $2b$'' concern the indicated rational coefficient maps. They do
not assert that the corresponding real or complex constants are
linearly independent over $\mathbb Q$. The source's existing formal
rank convention is preserved.

\item In the multiple-mark formula differentiate independent parameters
before specialization. A total derivative along the diagonal is a
different expression. For unequal marks, division of the symmetric
sum by $k!$ is not valid. For identical marked pairs it is valid.

\item The Gamma expressions use differences of analytic log-Gamma
branches on the right half-plane. An independently chosen logarithm
of a quotient need not have the same branch. The Stieltjes primitive
vanishes at its stated lower endpoint.

\item In the spectral decomposition the zero-jet correction
$\mathcal L_W$ is only the contribution of $F_0(B)$. Lower derivatives
$F_1,F_2,\ldots$ can still contain the Stieltjes and residue terms
present in the existing one-mass formulas. Removing $\mathcal L_W$
does not remove all such terms. The pair term $\zeta(2)$ in
\eqref{sd:eq:fourth} is required by the Gamma-ratio normalization.

\item A simplex face term is not an irreducible connected cumulant.
It can be nonzero when only two of several masses coincide. Formula
\eqref{sd:eq:quadraticface} gives its precise quadratic behavior and
prevents an incorrect stronger collision claim.
\end{enumerate}

An independent audit strengthened the proof of the polylogarithm
generator by replacing generic endpoint-regularization language with
an explicit convergent-MZV formula for the sole leading-one endpoint.
The Laurent derivative estimate was also written with its required
logarithmic factor. These are corrections made to the present draft
during review, not errors attributed to the baseline.

\subsection{What the reproducible checks establish}

The package separates proof certificates from numerical diagnostics.
The essential exact checks are as follows.

\begin{itemize}
\item A finite Faulhaber calculation, independent of the new Bernoulli
recurrence, agrees with it in 220 rational checks across powers one
through eight, including nonpositive anchors and boundary parity.
It retains MZVs as formal words. A separate finite-cutoff replay proves
the reduction of \eqref{tp:eq:main5} using only four classical MZV
relations. Rational coefficient extraction reproduces the full
eight-by-eight and eleven-by-eleven determinants and elimination rows.
For the repeated even-order extension, an independent standard-library
program checks 16 cubic cutoff polynomial identities and 192 finite
tail-difference versus strict-composition anchor identities.

\item The order-ten replay reconstructs the inherited determinant
polynomial from 120 signed permutation products. It checks all
$10{,}403$ integer monomials, the Bernoulli modular substitution,
and the small base matrices. A separate program computes all 25
coefficient norms used in the eventual argument and verifies the
integer factorial comparison. The new exceptional certificate uses
the displayed five-by-five matrix and independently computed Bernoulli
residues. No numerical determinant decides nonvanishing.

\item The marked-gap replay checks the formal insertion and partition
algebra through five independent labels. At five labels there are 52
set partitions and 120 surviving permutations; all unwanted endpoint
factors and shorter words cancel. The analytic insertion theorem is
proved in the text; the program checks its finite combinatorial
consequences.
\end{itemize}

Independent numerical representations then test the conclusions.
Three ordinary fifth- and sixth-power identities agree to more than
90 decimal digits using finite direct sums and centered Hurwitz tail
acceleration at two different cutoffs. The noninteger-endpoint marked
gap and its actual zero-marker derivative are reconstructed from
ordinary nested sums and compared with the decorated formulas; the
latter residual at cutoff 160 is approximately $3.73\times10^{-31}$.
Finite-interval Gamma and Stieltjes specializations agree to the stated
working precision. For the three-mass spectral third derivative, a
direct spectral evaluation agrees with independent Gamma face
quadrature to approximately $8.2\times10^{-32}$ at the highest
reported quadrature order.

These errors are observed cross-representation residuals. They are not
interval enclosures, and common truncation errors can survive a
comparison. The theorems rest on the analytic proofs and finite exact
identities. The package records the working precision, cutoff, method,
and exact-versus-numerical status of the checks. Its scripts use local
files only and need no repository checkout or network access to run.

The archive includes a standalone \LaTeX{} source, the compiled PDF,
modular section sources, Python verification programs, successful output
records, a claim-status ledger, source hashes, integration notes, and a
manifest verifier. The standalone source is generated from the modular
article, so the readable PDF and reusable section fragments share the
same mathematics. The source repository is left unchanged.

\section{Further research questions and exact next targets}
\label{sec:research}

The most useful next questions ask for identities, finite reductions,
or proofs of specified families. The following targets arise directly
from the present theorems.

\begin{enumerate}[label=\textbf{R\arabic*.}]
\item \textbf{Determine the smallest formal coordinate system at every
tail power.}
The universal generator supplies a finite span and exact ranks eight
and eleven for the chosen fifth- and sixth-power systems. Construct
canonical reduced coordinates at power $d$, state the admitted MZV
relations, and prove the rank of the resulting moment coefficient map.
One should distinguish the rank of that formal map from the unknown
arithmetic dimension after evaluating its coordinates.

\item \textbf{Classify compact polynomial annihilators.}
Equations \eqref{tp:eq:main5} and \eqref{tp:eq:main6} cancel all new
high-weight anchors in selected moment combinations. Find a structural
description of the polynomial weights producing a prescribed lower
coordinate span. A useful theorem would give the least degree and a
closed basis of such weights, with their exact asymptotic subtraction
polynomials. This could explain why the relatively small coefficients
in \eqref{intro:quintic} exist.

\item \textbf{Reduce the new anchored generators to explicit
polylogarithmic forms.}
The generator theorem is a finite algorithm with a depth upper bound.
Write a compact explicit expression for $H_5$ or $H_6$, then determine
which identities among allowed arguments and products reduce it.
Theorem statements should specify the function class and branches.
A proof of minimal depth would require an appropriate functional
invariant and is not supplied by the finite coordinate determinants.

\item \textbf{Find arithmetic recurrences for the resonant Laurent
family.}
Theorem~\ref{tp:thm:Laurent} fixes every coefficient and its spectral
germ. Determine whether the dependence on $r$ can be described by an
anchored recurrence with a controlled set of new logarithmic constants.
The first concrete case is the quartic constant term
$\Lambda_{4,r,0}$. Any proposed recurrence must reproduce the
$h_{4,r}\gamma_0(1/2)$ contribution exactly.

\item \textbf{Complete smaller order twelve, then remove the fixed-order
restriction.}
The order-ten proof shows how a modular certificate and a proved
eventual range can complement each other. Carry out the six-dimensional
contraction for smaller order twelve, account for every exceptional
residue below a sufficient threshold, and provide finite exact
certificates. A uniform determinant identity preventing cancellation
for all unequal even orders would settle the broader optimal-prefix
conjecture. The full gap-six family is still an intermediate target.

\item \textbf{Find closed reconstruction rows.}
The ordinary identities \eqref{mt10:ordinary} are explicit finite
rational algorithms involving an inverse matrix. Seek a fixed-size
kernel or product formula for its rows, including all endpoint
constants. This is a stronger goal than nonsingularity and may yield
short identities at large orders without constructing a large inverse.

\item \textbf{Resolve the antisymmetric two-marker component.}
For unequal marked pairs, the partition theorem gives the symmetric
combination. Determine an independent expression for
$\mathcal H(s,u,v,t)-\mathcal H(s,v,u,t)$ and its mixed derivatives.
A proposed reduction should retain the given relative quotient at
noninteger endpoints and prove an ordinary convergence domain before
continuation. This is the precise part of the several-marker question
not supplied by a classical symmetric-sum law.

\item \textbf{Combine marked clusters with nonpositive deletions.}
Unmarked deletions can shift distinct marked centers differently.
Classify exactly when such operations preserve the identical-cluster
or symmetric span, and identify the additional ordered coordinates
when they do not. The natural first test is one nonpositive integer
letter adjacent to a two-marker cluster. Independent differentiation
must precede every diagonal specialization.

\item \textbf{Evaluate special Gamma face integrals.}
For rational shifts and selected algebraic masses, seek identities
expressing $I_S(F_1)$ in a specified polylogarithmic, multiple-Gamma,
or zeta-derivative class. The quadratic identity
\eqref{sd:eq:quadraticface} gives a normalization check for every face
size. A finite exact integral is already available at arbitrary masses;
a special-function reduction would be an additional theorem.

\item \textbf{Give higher collision coefficients face by face.}
The quadratic coefficient in \eqref{sd:eq:collisionleading} is known
for every $k$. Derive closed symmetric-polynomial expressions for
$I_S(B^h)$ for general $h$, and for the logarithmically weighted
integrals in the higher derivative formula. Determine a generating
function that remains regular when any subset of masses coalesces.
This asks for explicit identities rather than only estimates.

\item \textbf{Allow independent spectral orders and arithmetic twists.}
Replace the common $s$ by independent $s_i$ in the product spectrum.
Derive the face decomposition for the resulting nonsymmetric Dirichlet
parameters, then compute prescribed transverse derivatives at their
joint origin. For periodic weights, project the resulting Gamma and
Stieltjes formulas onto Dirichlet characters while retaining the full
logarithmic normalization in the period. The mean-zero cancellation
provides an exact test at every polynomial weight.

\item \textbf{Return to the Gaussian period targets through an explicit
map.}
The new higher-depth polylogarithmic generators may supply functional
identities outside previously explored relation spaces. Before claiming
progress on $S_6$ or revised $S_8$, construct a precise map from one
such identity to the canonical residual, including its lower-weight
products and branch constants. A formal obstruction in a chosen
relation space and a high-precision small residual answer different
questions; neither alone proves or disproves the period identity.

\item \textbf{Formalize the finite-to-analytic interfaces.}
The Bernoulli boundary polynomial, finite summation identities,
modular determinant certificate, set-partition inversion, and face
M\"obius inversion are suitable finite formalization targets.
Separately formalize normal convergence of the subtracted spectral
families and agreement on their defining half-planes. Validated
interval quadrature would strengthen the computational diagnostics,
but remains distinct from formal proof of the exact identities.
\end{enumerate}

% END INCLUDED FILE: sections/06_audit_and_research.tex

\clearpage
\begin{thebibliography}{99}
\raggedright

% BEGIN INCLUDED FILE: references.tex
\bibitem{RepoMain}
Vladimir Reshetnikov and contributors,
\emph{ProveIt: Analysis/Polylogarithms, unified manuscript and research reports},
repository snapshot \texttt{da2792fc9cc8ef19c6a314be6c65e1437fbdc66b},
inspected 11 October 2026.
\url{https://github.com/VladimirReshetnikov/ProveIt}.

\bibitem{RepoQuartic}
\emph{Harmonic-Tail Reconstruction, Quartic Trigamma Moments,
and Spectral Stieltjes Calculus},
research continuation prepared for Vladimir Reshetnikov, 11 October 2026.
Incoming archive \path{ProveIt_Quartic_Spectral_Reconstruction_2026-10-11.zip}
in the source snapshot of \cite{RepoMain}.

\bibitem{RepoAlternating}
\emph{Alternating Tails and Spectral Defects},
research continuation prepared for Vladimir Reshetnikov, 11 October 2026.
Incoming archive \path{ProveIt_Alternating_Tails_Spectral_Defects_2026-10-11.zip}
in the source snapshot of \cite{RepoMain}.

\bibitem{DLMF}
F. W. J. Olver and others, eds.,
\emph{NIST Digital Library of Mathematical Functions}.
Sections 24.4 (Bernoulli polynomial identities) and 25.11 (Hurwitz zeta),
accessed 11 October 2026.
\url{https://dlmf.nist.gov/24.4}; \url{https://dlmf.nist.gov/25.11}.

\bibitem{HoffmanTails}
M. E. Hoffman,
Sums of products of Riemann zeta tails,
\emph{Mediterranean Journal of Mathematics} \textbf{13} (2016), 2771--2781.
\url{https://doi.org/10.1007/s00009-015-0653-9}.
\url{https://arxiv.org/abs/1602.03399}.

\bibitem{gm:INY}
K.~Ihara, Y.~Nakamura, and S.~Yamamoto,
\emph{Interpolant of truncated multiple zeta functions},
The Ramanujan Journal \textbf{67}, article~1 (2025).
\href{https://doi.org/10.1007/s11139-025-01045-2}{doi:10.1007/s11139-025-01045-2}.
See Corollary~4.2 and Theorem~5.2; the convolution representation is
Proposition~5.3.


\bibitem{gm:Hoffman1992}
M.~E.~Hoffman, \emph{Multiple harmonic series},
Pacific Journal of Mathematics \textbf{152} (1992), 275--290.
\href{https://msp.org/pjm/1992/152-2/pjm-v152-n2-p04-s.pdf}{Publisher PDF}.
See Theorem~2.2.


\bibitem{gm:Hoffman2018}
M.~E.~Hoffman, \emph{Quasi-shuffle algebras and applications},
arXiv:1805.12464 (2018).
\href{https://arxiv.org/abs/1805.12464}{arXiv:1805.12464}.
See Theorems~12--13 for the symmetric-sum identity and its general
quasi-shuffle form.


\bibitem{gm:gaps-source}
ProveIt research continuation,
\emph{Harmonic Gaps and Mixed Moments}, 11 October 2026.
Baseline commit \texttt{da2792fc9cc8ef19c6a314be6c65e1437fbdc66b};
\texttt{Analysis/Polylogarithms/docs/reports/} continuation
\texttt{stieltjes-harmonic-resonance/continuations/}
\texttt{harmonic-gaps-mixed-moments/}.
The one-marker construction is in Section~3; the present targets are
Q4 and Q5 in Section~6.

\bibitem{sd:HigherDefects}
\emph{Beta Polarization of Higher Spectral Defects:
All-order Stieltjes counterterms, Gamma integrals,
polylogarithmic second jets, and confluent quadratic masses},
research continuation prepared for Vladimir Reshetnikov, 11 October 2026.
Incoming archive \texttt{ProveIt\_Higher\_Spectral\_Defects\_2026-10-11.zip},
in the snapshot specified in \cite{RepoMain}.


\bibitem{sd:Carlson1987}
B. C. Carlson,
Dirichlet averages of $x^t\log x$,
\emph{SIAM Journal on Mathematical Analysis} \textbf{18} (1987), 550--565.
\url{https://doi.org/10.1137/0518043}.


\bibitem{sd:Carlson1991}
B. C. Carlson,
B-splines, hypergeometric functions, and Dirichlet averages,
\emph{Journal of Approximation Theory} \textbf{67} (1991), 311--325.
\url{https://doi.org/10.1016/0021-9045(91)90006-V}.

% END INCLUDED FILE: references.tex

\end{thebibliography}
\end{document}
